% Self-contained source generated by build.py.

% BEGIN preamble.tex
\documentclass[11pt,a4paper]{article}
\usepackage[T1]{fontenc}
\usepackage[utf8]{inputenc}
\usepackage{lmodern}
\usepackage[margin=26mm,headheight=14pt]{geometry}
\usepackage{amsmath,amssymb,amsthm,mathtools}
\usepackage{mathrsfs}
\usepackage{microtype,booktabs,longtable,array,enumitem,xcolor}
\usepackage{fancyhdr}
\usepackage{hyperref}
\hypersetup{colorlinks=true,linkcolor=blue!45!black,citecolor=blue!45!black,
 urlcolor=blue!55!black,pdftitle={Exact Resonant Identities Beyond the Diagonal},
 pdfauthor={Research continuation prepared with ChatGPT for the ProveIt project}}
\usepackage{bookmark}
\numberwithin{equation}{section}
\newtheorem{theorem}{Theorem}[section]
\newtheorem{lemma}[theorem]{Lemma}
\newtheorem{proposition}[theorem]{Proposition}
\newtheorem{corollary}[theorem]{Corollary}
\theoremstyle{definition}
\newtheorem{definition}[theorem]{Definition}
\theoremstyle{remark}
\newtheorem{remark}[theorem]{Remark}
\providecommand{\shuffle}{\mathbin{\amalg}}
\newcommand{\FP}{\operatorname{FP}}
\newcommand{\CT}{\operatorname{CT}}
\newcommand{\src}[1]{\nolinkurl{#1}}
\setlist{itemsep=3pt,topsep=5pt}
\setlength{\emergencystretch}{2em}
\allowdisplaybreaks[1]
\pagestyle{fancy}
\fancyhf{}
\fancyhead[L]{\small Exact resonant identities}
\fancyhead[R]{\small ProveIt research continuation}
\fancyfoot[C]{\thepage}
\renewcommand{\headrulewidth}{0.3pt}
\title{\textbf{Exact Resonant Identities Beyond the Diagonal}\\[5pt]
 \large Ordered Hurwitz jets, coincident Stieltjes products,\\
 harmonic powers, and a complete Cayley obstruction}
\author{Research continuation prepared with ChatGPT\\
 for Vladimir Reshetnikov's ProveIt project}
\date{11 October 2026}

% END preamble.tex

\begin{document}
\maketitle

% BEGIN sections/00_abstract.tex
\begin{abstract}
We continue the exact-identity programme of the ProveIt manuscript and
its incoming research reports at a fixed repository snapshot. Four
directions yield proved results. First, we determine the complete
Laurent expansion through the finite part of the depth-four ordered
Hurwitz zeta function along every transverse straight direction at
$(1,1,1,1)$. Three additional coordinates beyond the previously known
ordering function suffice; we give convergent representations, exact
shift laws, a differential relation, a three-direction reconstruction,
and the cancellation locus of the displayed nonsymmetric terms.
Second, a coordinate-completed cotangent--Hurwitz transform evaluates
one arbitrary Stieltjes derivative times any finite product of
coincident reflected polygamma factors. Its resonant subtraction retains
the whole parameter-dependent amplitude. Consequences include cubic
differences, higher reflected moments, and log-Gamma families involving
zeta derivatives and harmonic numbers. Third, the Dirichlet series of
raw harmonic powers with a continuously shifted outer denominator has
polynomial Laurent coefficients and finite parts at every nonpositive
integer. Centering the shift at $1/2$ removes every nonpositive even
pole, and the resulting values belong to the ordinary multiple-zeta
algebra without requiring Euler's constant. Fourth, an exact integer
functional extends the existing bounded weight-seven obstruction to
the entire convergent Cayley shuffle ideal plus the specified additional
finite relation family. This constrains a proof strategy for $S_6$;
it does not decide the period conjecture.

We provide complete analytic and algebraic arguments, a corrected sign
in a specifically identified harmonic-power preprint formula, explicit
integration guidance, and reproducible exact and numerical checks.
The numerical diagnostics are distinguished throughout from proofs and
interval certificates. The short $S_6$, revised $S_8$, generic cubic
Tornheim reduction, and arithmetic minimality questions remain open.
\end{abstract}
\noindent\textbf{Keywords.} Multiple polylogarithm; ordered Hurwitz zeta;
generalized Stieltjes constant; Hadamard finite part; polygamma reflection;
harmonic-power Dirichlet series; Euler sums; shuffle algebra.

% END sections/00_abstract.tex

\clearpage
{\small\tableofcontents}
\clearpage

% BEGIN sections/01_scope.tex
\section{Scope, antecedents, and the exact meaning of the results}
\label{scope:section}

The starting point is the canonical ProveIt manuscript
\cite{RepoMain} and the eleven archives present in \src{docs/incoming}
at repository commit
\begin{center}
\texttt{5a790187c8e186e41e2b990b4941cb7a1a3c7b6b}.
\end{center}
The commit timestamp is 11 October 2026, 00:09:09 UTC. The incoming
reports retain their 10 October dates. The delivered source manifest
records the retrieved paths and hashes. This article is an integration
candidate; it does not modify the canonical manuscript or promote
unreviewed incoming statements by editorial fiat.

The immediate motivation is three explicit unfinished parts of the
latest exact-identity programme. The ordered-Hurwitz report
\cite{RepoOrdered} asks for the first independent-slope formula beyond
depth three and for the directions where its additional ordering
coefficients disappear. The triple-Stieltjes report \cite{RepoTriple}
provides finite closure for reflected factors with distinct seams, while
coincidences require a separate normalization. The Gaussian depth
projector report \cite{gf:Depth} proposes completing the Cayley ideal
before repeating a finite reduction search. Sections~\ref{oq:section},
\ref{ct:section}, and~\ref{gf:sec} answer these questions in the
precise senses stated below. The harmonic-power work in
Section~\ref{hp:sec:main} gives a further exact family connecting
multiple polylogarithms, Gamma coefficients, harmonic numbers, and
spectral finite parts.

\subsection{A ledger of contributions}

\begin{longtable}{@{}p{0.22\textwidth}p{0.39\textwidth}p{0.32\textwidth}@{}}
\toprule
Direction & Proved here & Boundary of the claim\\
\midrule
\endhead
Depth-four ordered resonance & A spanning normal form for all transverse
straight rays, convergent coordinates, shifts, and reconstruction &
No arithmetic independence; no prescription for a ray contained in a
polar hyperplane.\\[5pt]
Coincident reflected products & All Stieltjes and argument derivative
indices; arbitrary finite reflected products; explicit Gamma primitives &
The generic unreflected cubic Tornheim coordinate is not evaluated.\\[5pt]
Raw harmonic powers & All nonpositive principal and finite coefficients
as shift polynomials; centered cancellation; explicit families &
The first regular spectral coefficient generally needs global Mellin
data.\\[5pt]
Gaussian formal obstruction & Nonmembership after adding the complete
Cayley ideal to the frozen additional relation span &
$S_6$ and revised $S_8$ remain conjectural as numerical period identities.\\[5pt]
Literature correction & A sign correction to equation (22) of
arXiv:2507.03058v1, with an all-logarithmic proof &
It is a correction to that identified formula, not a rejection of the
preprint's entire framework.\\
\bottomrule
\end{longtable}

The word ``new'' in this report means an explicitly demonstrated
continuation beyond the inspected ProveIt results. We identify the
classical inputs and the nearest primary-source antecedents. We do not
claim an exhaustive priority search over every formulation in the
literature. In particular, the Gamma diagonal of ordered regular germs,
ordinary Hurwitz Fourier expansions, shuffle regularization, and
raw-power meromorphic continuation are antecedents.

The analytic context is the Laurent theory of multiple zeta functions
developed by Saha \cite{Saha} and by Matsumoto, Onozuka, and Wakabayashi
\cite{MOW}. Hurwitz integral methods and their Gamma applications have
substantial classical development, including Espinosa and Moll
\cite{EspinosaMoll1,EspinosaMoll2}. Tin \cite{hp:Tin2025} studies raw
harmonic-power Dirichlet series and their harmonic Stieltjes constants;
Karg\i n, Dil, Cenkci, and Can \cite{Kargin} study related harmonic zeta
constants. We retain their strict/inclusive and inner/outer shift
distinctions rather than silently identifying different series.
The shuffle-polynomial and Eulerian-logarithm structure in the formal
argument is classical \cite{gf:Radford,gf:Patras}; the new certificate
concerns its application to the complete ideal intersection.

\subsection{Three normalizations that must not be interchanged}

For $a>0$, throughout the article,
\[
 \zeta(1+u,a)=\frac1u+
       \sum_{m\ge0}\frac{(-1)^m}{m!}\gamma_m(a)u^m,
 \qquad \gamma_0(a)=-\psi(a).
\]
A spectral derivative differentiates the order $s$ of a zeta function;
an argument derivative differentiates its positive shift $a$ or its
position $x$. A superscript $(p)$ on $\gamma_m$ or $\psi$ in the
reflected section always denotes an argument derivative. Bracketed
superscripts $[j]$ on $z_k$ in the ordered section denote spectral
derivatives. The ordinary constants are $\gamma_m=\gamma_m(1)$.
Generalized harmonic numbers are
$H_n^{(r)}=\sum_{j=1}^n j^{-r}$, with $H_0^{(r)}=0$.

There are three different finite-coefficient operations.
\begin{enumerate}
\item A \emph{spectral finite part} is the constant Laurent coefficient
in a specified complex parameter. The ray variable $\varepsilon$ in
Section~\ref{oq:section} and the spectral variable $s+m$ in
Section~\ref{hp:sec:main} are explicitly fixed.
\item A \emph{coordinate finite-part integral} is the constant term of
a cut-off integral in fixed local coordinates. On $(0,1)$ those
coordinates are $x$ and $1-x$. Powers and logarithms are expanded before
the constant term is selected. This is the convention in
Section~\ref{ct:section}.
\item A \emph{formal shuffle relation} is an equality in a specified
rational word algebra. Its truth and exact verification are independent
of an unproved assertion that the period map is injective.
\end{enumerate}
All ordinary integrals, convergent sums, analytic continuations, and
finite parts are identified where they first occur. Analytic continuation
of an integral and coefficientwise coordinate completion agree away
from a resonance, but can differ by a finite Taylor correction at one.
That correction is computed explicitly in the coincident theorem.

The multiple-polylogarithm convention used for the Gaussian target is
\[
 \operatorname{Li}_{r,s}(u,v)
 =\sum_{n>m\ge1}\frac{u^nv^m}{n^rm^s}.
\]
The one-variable multiple polylogarithms in the harmonic-power kernel
carry the argument only on the outer index. These conventions fix the
word order, colors, and harmonic conversions; changing them changes
intermediate coordinate formulas.

\subsection{Reading and integration map}

The four mathematical sections can be read independently after the
normalizations above. The ordered section supplies all of its required
lower-depth germs. The reflected section proves the completed transform
before using polynomial reduction. The harmonic section gives both a
finite general algorithm and closed low-weight examples. The formal
section includes the algebraic argument connecting its finite exact
checks to the unrestricted ideal. Section~\ref{audit:section} distinguishes
an actual correction from normalization safeguards, and
Section~\ref{research:section} identifies concrete remaining targets.

For canonical integration, the ordered section naturally follows the
latest ordered-Hurwitz contribution; the reflected section follows the
coincident/contact calculus; the harmonic section belongs with harmonic
polylogarithm and endpoint regularization material; and the complete
Cayley certificate belongs immediately after the existing bounded
$S_6$ obstruction. The retained proof-status qualifications are part of
each mathematical statement.

% END sections/01_scope.tex


% BEGIN sections/02_ordered.tex
\section{Independent directions at depth four}
\label{oq:section}

The most recent ordered-Hurwitz report asks for an explicit description of
the first independent-slope case beyond depth three
\cite[Questions 1--2]{RepoOrdered}. We give such a description. Three
additional coordinates suffice beyond its previously defined ordering
function $\eta(a)$. All three have absolutely convergent representatives.
The resulting formula includes the complete principal part, every
transverse straight direction, exact shift recurrences, and a finite
reconstruction of the coordinates from three positive directions.
The construction is a spanning normal form; it does not assert arithmetic
independence of its coordinates.

\subsection{Conventions and the holomorphic regular germs}

Fix $a>0$ and use real logarithms of $n+a$. Set
\[
 Z_d(s_1,\ldots,s_d;a)
 =\sum_{n_1>\cdots>n_d\ge0}\prod_{j=1}^d(n_j+a)^{-s_j},
 \qquad Z_0=1.
\]
The initial domain is $\Re(s_1+\cdots+s_j)>j$ for every $j$.
Write $u_j=s_j-1$ and $P_k=u_1+\cdots+u_k$.
The Laurent convention is
\begin{equation}\label{oq:stieltjes}
 \zeta(1+u,a)=\frac1u+
 \sum_{j\ge0}\frac{(-1)^j\gamma_j(a)}{j!}u^j.
\end{equation}
Throughout this section,
\[
 g_j=\gamma_j(a),\quad g=g_0=-\psi(a),\quad
 z_k^{[j]}=\left.\partial_s^j\zeta(s,a)\right|_{s=k},
 \quad z_k=z_k^{[0]}.
\]
Thus a prime on an ordinary zeta value denotes its order derivative,
whereas $\eta'(a)$ below denotes its shift derivative.

For completeness, the regular germs can be constructed directly.
Define
\begin{align}
 D(s,x)&=\zeta(s,x+1)-\frac{x^{1-s}}{s-1},\label{oq:D}\\
 E_d(\mathbf u;a)&=
 \sum_{n_2>\cdots>n_d\ge0}
 D(1+u_1,n_2+a)\prod_{j=2}^d(n_j+a)^{-1-u_j}.
 \label{oq:E}
\end{align}
Let $R_0=1$, $R_1(u;a)=\zeta(1+u,a)-1/u$, and recursively put
\begin{align}\label{oq:Rrec}
 R_d(\mathbf u;a)=E_d(\mathbf u;a)+
 \frac{R_{d-1}(u_1+u_2,u_3,\ldots;a)
       -R_{d-1}(u_2,u_3,\ldots;a)}{u_1}.
\end{align}
The divided difference takes its removable value when $u_1=0$.

\begin{lemma}[Ordered decomposition and normalization]\label{oq:regular}
The germs $R_d$ are holomorphic on $\sum_j|u_j|<1/2$, and
\begin{equation}\label{oq:decomp}
 Z_d(1+\mathbf u;a)=
 \sum_{k=0}^d\frac{R_{d-k}(u_{k+1},\ldots,u_d;a)}{P_1\cdots P_k}.
\end{equation}
Furthermore, if $B_d(a)=R_d(0,\ldots,0;a)$, then
\begin{equation}\label{oq:gamma}
 \sum_{d\ge0}B_d(a)y^d=\frac{\Gamma(a)}{\Gamma(a+y)}.
\end{equation}
\end{lemma}
\begin{proof}
For $\Re s>0$, the tail defect has the integral
\[
 D(s,x)=\frac1{\Gamma(s)}\int_0^\infty
 t^{s-1}e^{-xt}\left(\frac1{e^t-1}-\frac1t\right)dt.
\]
The bracket is bounded at zero. This proves removability at $s=1$
and $D(s,x)=O(x^{-\Re s})$ on compact subsets. The summands in
\eqref{oq:E}, after summing the inner finite indices, have a
uniform bound of the form
$C(1+n_2)^{-2+\sum|u_j|}(1+\log(1+n_2))^{d-2}$.
Normal convergence and the integral expression for a divided difference
prove holomorphy. Summing first over $n_1>n_2$ gives
\[
 Z_d(1+\mathbf u;a)=\frac1{u_1}
 Z_{d-1}(1+u_1+u_2,1+u_3,\ldots;a)+E_d(\mathbf u;a),
\]
so induction proves \eqref{oq:decomp}. On the common diagonal,
the elementary-symmetric product for $Z_d$ implies, coefficientwise
as a formal power series in $y$,
\begin{align}\label{oq:diaggenerator}
 \sum_{d\ge0}R_d(t,\ldots,t;a)y^d
 =\exp\left\{yR_1(t;a)+
 \sum_{m\ge2}\frac{(-1)^{m-1}}m\zeta(m+mt,a)y^m\right\}.
\end{align}
At $t=0$, the logarithmic Taylor expansion of $\Gamma(a+y)$
gives \eqref{oq:gamma}.
\end{proof}

This construction and the diagonal generating function are antecedents,
not new claims of this article. For $a=1$, the ordered Laurent framework
is due to Saha \cite{Saha}; general integer-point Laurent expansions are
studied in \cite{MOW}. The common-shift construction above is the one
in \cite{RepoOrdered}. In particular,
\begin{align}
 B_2&=\frac{g^2-z_2}{2},&
 B_3&=\frac{g^3-3gz_2+2z_3}{6},\label{oq:B23}\\
 B_4&=\frac{g^4-6g^2z_2+3z_2^2+8gz_3-6z_4}{24}.
 \label{oq:B4}
\end{align}

\subsection{The small jet that controls the answer}

Define
\begin{align}
 \eta(a)&=[u]R_2(u,0;a),\label{oq:eta}\\
 \rho(a)&=[u^2]R_2(u,0;a)-[v^2]R_2(0,v;a),\label{oq:rho}\\
 \kappa(a)&=(\partial_1-2\partial_2+\partial_3)R_3(0,0,0;a),
 \label{oq:kappa}\\
 T(a)&=(\partial_u-\partial_v)Z_2(2+u,1+v;a)\big|_{u=v=0}.
 \label{oq:T}
\end{align}
The factorial convention in \eqref{oq:rho} matters: $\rho$ is a
difference of Taylor coefficients, not a difference of second derivatives.
The last coordinate has the elementary convergent series
\begin{equation}\label{oq:Tseries}
 T(a)=\sum_{n>m\ge0}
 \frac{\log(m+a)-\log(n+a)}{(n+a)^2(m+a)}.
\end{equation}
Introduce the ordinary-data combinations
\begin{align}
 D_1&=-gg_1-z_2^{[1]},\label{oq:D1}\\
 E&=\frac{gg_2-z_2^{[2]}}2,
 &F&=\frac{g_1^2-z_2^{[2]}}2,\label{oq:EF}\\
 S&=-g_1B_2-gz_2^{[1]}+z_3^{[1]},\label{oq:S}\\
 V&=g\eta+\frac{g_1(g^2+z_2)}2+gz_2^{[1]}-T.
 \label{oq:V}
\end{align}
Only $V$ retains a nonordinary coordinate; the letters $D_1,E,F,S$
denote explicitly displayed finite expressions.

\begin{theorem}[Complete required jet]\label{oq:jet}
The quadratic Taylor polynomial of $R_2$ and the linear Taylor polynomial
of $R_3$ are
\begin{align}\label{oq:R2jet}
 R_2(u,v;a)={}&B_2+\eta u+(D_1-\eta)v
 +\frac{E+\rho}{2}u^2+Fuv+\frac{E-\rho}{2}v^2
 +O((|u|+|v|)^3),\\
 R_3(u,v,w;a)={}&B_3+
 \frac S3(u+v+w)+\frac\kappa6(u-2v+w)
 +\frac V2(u-w)+O((|u|+|v|+|w|)^2).
 \label{oq:R3jet}
\end{align}
\end{theorem}
\begin{proof}
The regularized depth-two stuffle identity is
\begin{equation}\label{oq:stuffle2}
 R_2(u,v)+R_2(v,u)=R_1(u)R_1(v)-\zeta(2+u+v,a).
\end{equation}
It follows directly by substituting \eqref{oq:decomp} into the ordinary
stuffle identity; the rational polar terms cancel. Its coefficients of
$u$, $u^2$ and $uv$, together with the definition of $\rho$, prove
\eqref{oq:R2jet}.

For the next depth put
\[
 C(u,v)=Z_2(2+u,1+v;a),\qquad
 L(u,v)=Z_2(1+u,2+v;a)-\frac{\zeta(2+v,a)}u.
\]
The first function converges normally near zero. The second is holomorphic
there because the two-factor stuffle law gives the exact expression
\begin{equation}\label{oq:L}
 L(u,v)=R_1(u)\zeta(2+v,a)-C(v,u)-\zeta(3+u+v,a).
\end{equation}
Applying the three-factor stuffle identity and cancelling the explicit
prefix poles gives
\begin{align}\label{oq:stuffle3}
 R_1(u)R_2(v,w)={}&R_3(u,v,w)+R_3(v,u,w)+R_3(v,w,u)
 +C(u+v,w)+L(v,u+w).
\end{align}
For clarity, the needed cancellation uses
$R_2(u,w)+R_2(w,u)=R_1(u)R_1(w)-\zeta(2+u+w,a)$.
Thus \eqref{oq:stuffle3} follows from meromorphic identities and entails
no unproved assumption that arbitrary finite-part prescriptions preserve
products.

Write $X=\partial_1R_3(0)$, $Y=\partial_2R_3(0)$ and
$Z=\partial_3R_3(0)$. Differentiating \eqref{oq:stuffle3} with respect
to $u$ and using \eqref{oq:L} gives
\[
 X+Y+Z=-g_1B_2-gz_2^{[1]}+z_3^{[1]}=S.
\]
Differentiating with respect to $v$ instead gives
\[
 2X+Y=g\eta+g_1z_2+z_3^{[1]}-T.
\]
Subtracting the preceding equation proves $X-Z=V$.
Finally $X-2Y+Z=\kappa$ by definition. Solving these three linear
equations gives \eqref{oq:R3jet}.
\end{proof}

\subsection{Every transverse depth-four Laurent constant}

For $p,b,c,d\in\mathbb C$ suppose
\begin{equation}\label{oq:transverse}
 p(p+b)(p+b+c)(p+b+c+d)\ne0.
\end{equation}
Here $\operatorname{FP}_{\varepsilon=0}$ means the constant coefficient
of the meromorphic function after the specified substitution. It is not
an iterated residue and is not a value on a polar hyperplane.

\begin{theorem}[Depth-four directional normal form]\label{oq:main}
Under \eqref{oq:transverse},
\begin{align}\label{oq:FP4}
 &\operatorname{FP}_{\varepsilon=0}
 Z_4(1+p\varepsilon,1+b\varepsilon,1+c\varepsilon,1+d\varepsilon;a)
 \notag\\
 &\quad=B_4+
 \frac1p\left\{\frac{S(b+c+d)}3+
            \frac{\kappa(b-2c+d)}6+\frac{V(b-d)}2\right\}
 \notag\\
 &\qquad\quad+
 \frac{E(c^2+d^2)/2+Fcd+\rho(c^2-d^2)/2}{p(p+b)}
 -\frac{g_3d^3}{6p(p+b)(p+b+c)}.
\end{align}
The complete principal part is
\begin{align}\label{oq:principal}
 &\frac{1}{p(p+b)(p+b+c)(p+b+c+d)\varepsilon^4}
 +\frac{g}{p(p+b)(p+b+c)\varepsilon^3}
 \notag\\
 &\quad+\frac1{\varepsilon^2}
 \left\{\frac{B_2}{p(p+b)}-
 \frac{dg_1}{p(p+b)(p+b+c)}\right\}
 \notag\\
 &\quad+\frac1\varepsilon
 \left\{\frac{B_3}p+
 \frac{\eta(c-d)+D_1d}{p(p+b)}+
 \frac{g_2d^2}{2p(p+b)(p+b+c)}\right\}.
\end{align}
\end{theorem}
\begin{proof}
At depth four, \eqref{oq:decomp} has five terms:
\[
 \frac1{P_1P_2P_3P_4}+
 \frac{R_1(u_4)}{P_1P_2P_3}+
 \frac{R_2(u_3,u_4)}{P_1P_2}+
 \frac{R_3(u_2,u_3,u_4)}{P_1}+R_4(\mathbf u).
\]
The constant term requires the cubic coefficient of $R_1$, the quadratic
coefficient of $R_2$, the linear coefficient of $R_3$, and $R_4(0)=B_4$.
Equations \eqref{oq:stieltjes}, \eqref{oq:R2jet} and \eqref{oq:R3jet}
give exactly \eqref{oq:FP4}. The lower powers give
\eqref{oq:principal}. Holomorphy of the remainders makes the coefficient
calculation legitimate, including directions with some zero or negative
slopes, provided the prefix sums remain nonzero.
\end{proof}

The answer depends on $p$ and the three trailing slopes, but it needs
no new fourth-depth regular constant. The new information lies in
lower-depth derivatives, exactly as the ordered pole geometry predicts.

\begin{corollary}[Exact cancellation locus in this normal form]
\label{oq:locus}
Treat $\rho,\kappa,T,\eta$ as named coordinates, without assigning
unproved relations among their values. Under \eqref{oq:transverse},
all their coefficients in \eqref{oq:FP4} vanish if and only if
$b=c=d$. On this locus the answer reduces to ordinary Hurwitz jets.
\end{corollary}
\begin{proof}
The coefficient of $T$ is $-(b-d)/(2p)$ and the coefficient of
$\kappa$ is $(b-2c+d)/(6p)$. Their simultaneous vanishing forces
$b=d$ and then $c=b$. This also kills the $\rho$ and $\eta$
coefficients. Conversely, substitution proves cancellation.
\end{proof}

This corollary classifies cancellation in the proved coordinate expression.
It does not say that no other direction can simplify after a new functional
identity is discovered. In particular, it is not a statement of arithmetic
independence at an individual value of $a$.

For $b=c=d=q$, the formula becomes
\begin{equation}\label{oq:diagonalcheck}
 B_4+\frac qpS+
 \frac{q^2}{p(p+q)}\left\{\frac{g_1^2+gg_2}{2}-z_2^{[2]}\right\}
 -\frac{q^3g_3}{6p(p+q)(p+2q)},
\end{equation}
which recovers the already proved distinguished-slope formula of
\cite{RepoOrdered}. This is a consistency specialization, not an additional
new evaluation.

\subsection{Three positive rays recover all three coordinates}

Define the ordinary part of \eqref{oq:FP4} by
\begin{align}\label{oq:ordinary}
 \mathcal O(p,b,c,d;a)=B_4+
 \frac{S(b+c+d)}{3p}+
 \frac{E(c^2+d^2)/2+Fcd}{p(p+b)}
 -\frac{g_3d^3}{6p(p+b)(p+b+c)}.
\end{align}
Let $\Delta_A,\Delta_B,\Delta_C$ be the directional finite parts minus
$\mathcal O$ at the respective rays
\[
 A=(1,1,2,1),\qquad B=(1,1,3,1),\qquad C=(1,2,1,1).
\]
All three lie in the positive convergence directions.

\begin{corollary}[Finite reconstruction]\label{oq:reconstruction}
One has
\begin{align}
 \rho&=2\Delta_B-4\Delta_A,\label{oq:recoverrho}\\
 \kappa&=\frac92\Delta_B-12\Delta_A,\label{oq:recoverkappa}\\
 V&=2\Delta_C-\frac32\Delta_B+4\Delta_A.\label{oq:recoverV}
\end{align}
Consequently $T$ is recovered from \eqref{oq:V} once $\eta$ is known.
\end{corollary}
\begin{proof}
The three equations supplied by \eqref{oq:FP4} are
\[
 \begin{pmatrix}\Delta_A\\\Delta_B\\\Delta_C\end{pmatrix}
 =\begin{pmatrix}-1/3&3/4&0\\-2/3&2&0\\1/6&0&1/2\end{pmatrix}
 \begin{pmatrix}\kappa\\\rho\\V\end{pmatrix}.
\]
The determinant is $-1/12$. Inverting this rational matrix gives the
displayed formulas.
\end{proof}

Thus the appearance of three coordinates is observable by a finite exact
linear calculation. The invertible matrix proves equivalence of two
spanning descriptions; it does not prove that the three numerical values
are independent over any field.

\subsection{Absolutely convergent representatives}

Put
\[
 Q_j(x)=\gamma_j(x+1)+\frac{\log^{j+1}x}{j+1},\qquad
 h_n(a)=\sum_{m=0}^{n-1}\frac1{m+a},\qquad
 \ell_n(a)=\sum_{m=0}^{n-1}\frac{\log(m+a)}{m+a}.
\]
Here empty sums vanish. The familiar identities
\[
 h_n(a)=\psi(n+a)-\psi(a),\qquad
 \ell_n(a)=\gamma_1(a)-\gamma_1(n+a)
\]
make the representatives numerically usable with one outer summation.

\begin{theorem}[Convergent coordinate formulas]\label{oq:series}
With $x=n+a$ in every sum,
\begin{align}
 \eta(a)&=\frac{g_2}{2}-\sum_{n\ge0}\frac{Q_1(x)}x,
 \label{oq:etaseries}\\
 \rho(a)&=\frac{g_3}{3}+
 \frac12\sum_{n\ge0}\frac{Q_2(x)-Q_0(x)\log^2x}{x},
 \label{oq:rhoseries}\\
 \kappa(a)&=\sum_{n\ge0}
 \frac{-Q_1(x)h_n(a)+Q_0(x)(2\log x\,h_n(a)-\ell_n(a))}{x}
 -\frac32(E+\rho)+F.\label{oq:kappaseries}
\end{align}
All these series, as well as \eqref{oq:Tseries}, converge absolutely
and locally uniformly for $a$ in compact subintervals of $(0,\infty)$.
\end{theorem}
\begin{proof}
Taylor expansion in \eqref{oq:D} gives
$D(1+u,x)=\sum_{j\ge0}(-1)^jQ_j(x)u^j/j!$.
If $r^{(d)}_{m_1,\ldots,m_d}=[\mathbf u^{\mathbf m}]R_d$, the
divided difference in \eqref{oq:Rrec} yields
\begin{equation}\label{oq:coeffrec}
 r^{(d)}_{m_1,\ldots,m_d}
 =[\mathbf u^{\mathbf m}]E_d+
 \binom{m_1+m_2+1}{m_1+1}
 r^{(d-1)}_{m_1+m_2+1,m_3,\ldots,m_d}.
\end{equation}
For $(m_1,m_2)=(1,0),(2,0),(0,2)$ this gives the first two formulas.
For $(1,0,0),(0,1,0),(0,0,1)$, their combination with coefficients
$1,-2,1$ gives the sum in \eqref{oq:kappaseries} and the correction
$-3[u^2]R_2+[uv]R_2=-3(E+\rho)/2+F$.
The estimates used for \eqref{oq:E}, followed by Cauchy's formula on
a slightly larger compact spectral polydisc, justify every coefficient
extraction and give absolute convergence. Alternatively,
$Q_j(x)=O((1+\log x)^j/x)$ suffices, since
$h_n=O(1+\log n)$ and $\ell_n=O(1+\log^2n)$.
\end{proof}

These formulas define the coordinates independently of a numerical
Laurent fit. They are useful even if no further ordinary-zeta reduction
exists, and they keep the unevaluated part confined to explicitly
convergent harmonic--Stieltjes sums.

\subsection{Exact shift and differential compatibility}

\begin{proposition}[Triangular shift laws]\label{oq:shift}
Let a superscript $+$ mean evaluation at $a+1$, and let $L=\log a$.
Then
\begin{align}
 \eta(a)-\eta(a+1)&=-\frac{g_1^+}{a},\label{oq:shifteta}\\
 \rho(a)-\rho(a+1)&=\frac{g_2^+-g^+L^2}{2a},\label{oq:shiftrho}\\
 T(a)-T(a+1)&=\frac{z_2^{[1],+}+Lz_2^+}{a},\label{oq:shiftT}\\
 \kappa(a)-\kappa(a+1)&=
 \frac{3\eta^+-2D_1^+-LB_2^+}{a}.\label{oq:shiftkappa}
\end{align}
\end{proposition}
\begin{proof}
The terms with last index zero give the exact ordered shift identity
\[
 Z_d(\mathbf s;a)-Z_d(\mathbf s;a+1)
 =a^{-s_d}Z_{d-1}(s_1,\ldots,s_{d-1};a+1).
\]
Substitution into \eqref{oq:decomp} and induction on depth show that
\begin{equation}\label{oq:Rshift}
 R_d(\mathbf u;a)-R_d(\mathbf u;a+1)
 =a^{-1-u_d}R_{d-1}(u_1,\ldots,u_{d-1};a+1).
\end{equation}
The first, second and fourth assertions are the indicated Taylor
coefficients of \eqref{oq:Rshift}. For the third, use
$C(u,v;a)-C(u,v;a+1)=a^{-1-v}\zeta(2+u,a+1)$ and apply
$\partial_u-\partial_v$.
\end{proof}

\begin{proposition}[A convergent derivative identity]\label{oq:derivative}
Let $C_0(a)=Z_2(2,1;a)$. Then
\begin{equation}\label{oq:etaT}
 T(a)=-\eta'(a)-C_0(a)+g_1z_2+z_3^{[1]}.
\end{equation}
At $a=1$, the classical Euler evaluation $Z_2(2,1;1)=\zeta(3)$ gives
\[
 T(1)=-\eta'(1)-\zeta(3)+\gamma_1\zeta(2)+\zeta'(3).
\]
\end{proposition}
\begin{proof}
Differentiating the ordered defining sum in its initial domain, and
then continuing, gives
\[
 \partial_aR_2(u,v;a)=-(1+u)C(u,v)-(1+v)L(u,v).
\]
The pole removed from $L$ cancels exactly with
$\partial_aR_1(v)/u=-(1+v)\zeta(2+v,a)/u$.
Apply $\partial_u$ at zero and substitute
$L_u(0,0)=-g_1z_2-z_3^{[1]}-C_v(0,0)$ from \eqref{oq:L}.
The result is \eqref{oq:etaT}. Normal convergence justifies the
derivatives of $C$; the regular germs are holomorphic in $a$ for
$\Re a>0$ on the principal branch.
\end{proof}

Equations \eqref{oq:shifteta}--\eqref{oq:etaT} connect the new directional
coordinates to generalized harmonic sums, generalized Stieltjes constants,
and parameter derivatives. They also supply independent consistency checks
for implementations. A complete reduction of $\rho$ or $\kappa$ to ordinary
single-zeta data is not proved here.

% END sections/02_ordered.tex


% BEGIN sections/03_reflected.tex
% Standalone section source for the parent article. Uses standard amsmath/amsthm.
\section{Coincident reflected Stieltjes products close in ordinary zeta jets}
\label{ct:section}

Let
\[
 K(x)=\pi\cot(\pi x),\qquad
 \zeta(1+u,x)-\frac1u
   =\sum_{m\ge0}\frac{(-1)^m}{m!}\gamma_m(x)u^m.
\]
In particular $\gamma_0(x)=-\psi(x)$. All finite parts in this section use
\emph{both} unit endpoint coordinates $x$ and $1-x$. Equivalently,
\[
 \operatorname{FP}\int_0^1 f(x)\,dx
 =\operatorname{CT}_{\epsilon\downarrow0}
       \int_\epsilon^{1-\epsilon} f(x)\,dx,
\]
where the constant coefficient is taken after expanding every endpoint
power and logarithm. Independent lower and upper cutoffs give the same
answer because the two local contributions are added. This convention
is invariant under $x\mapsto1-x$, and it assigns zero to
$\operatorname{FP}\int_0^1x^{-1}\log^j x\,dx$.

The incoming triple report proves finite closure for one Stieltjes factor
\cite{RepoTriple} and any number of reflected polygamma factors whose seams are pairwise
separated. The result below allows every reflected factor to collide with
the Stieltjes seam. It is an exact closed subclass of the coincident
product problem. It does not evaluate the unreflected cube
$\operatorname{FP}\int\psi^3$, whose remaining diagonal Tornheim third
jet is a different coordinate. The Fourier transform and reflection
inputs are classical; the contribution here is the full coordinate
completion and its all-order coincident reflected consequences.

\subsection{A classical meromorphic transform and its local completion}

\begin{lemma}[Cotangent--Hurwitz transform]
\label{ct:transform}
For $r\ge0$, initially when $\Re s$ is sufficiently negative and then by
meromorphic continuation,
\begin{equation}\label{ct:J}
 J_r(s):=\operatorname{FP}\int_0^1\zeta(s,x)K^{(r)}(x)\,dx
 =(s)_r\,\pi\cot\frac{\pi(s+r)}2\,\zeta(s+r).
\end{equation}
At a spectral resonance the right side denotes the meromorphic
continuation, not the coordinate finite part of a specialized
coefficient.
\end{lemma}
\begin{proof}
The periodic principal value of $K$ has nonzero Fourier coefficients
$-\pi i\operatorname{sgn}(n)$ and zero mean. This follows, for example,
by differentiating the Abel Fourier series for
$\log(2\sin\pi x)$ in distributions. The Hurwitz Fourier expansion \cite[Eq.~25.11.9]{DLMFHurwitz} gives,
for $\Re s<0$,
\[
 J_0(s)=2\pi\Gamma(1-s)(2\pi)^{s-1}
             \cos(\pi s/2)\zeta(1-s).
\]
The sum is absolutely convergent because its summands have magnitude
$O(n^{\Re s-1})$. The Riemann functional equation \cite{DLMFZetaReflection} turns the displayed
expression into $\pi\cot(\pi s/2)\zeta(s)$, with removable singularities
interpreted by continuation.

For $\Re s<-r-1$, the singular term $x^{-s}$ and sufficiently many
of its derivatives vanish at zero. The smooth part
$\zeta(s,1+x)$ has the same periodic endpoint jets as $\zeta(s,x)$ at
one. Consequently distributional integration by parts against
$K^{(r)}$ agrees with the stated endpoint finite part and gives
\[
 J_r(s)=(-1)^r\operatorname{FP}\int_0^1
              \partial_x^r\zeta(s,x)K(x)\,dx
       =(s)_rJ_0(s+r).
\]
There are no unrecorded endpoint constants in this initial region:
the smooth jets match and the remaining power vanishes. Finite endpoint
Taylor subtraction supplies meromorphic continuation in $s$, proving
\eqref{ct:J} everywhere as a meromorphic identity.
\end{proof}

Set $N=p+r$ and define $P_j(u)=(1+u)_j$, so $P_0=1$.
For $p,r\ge0$ put
\begin{align}
 \mathcal R_{p,r}(u)
 =(-1)^p\bigg[&P_N(u)\,
      \pi\cot\frac{\pi(N+1+u)}2\,\zeta(N+1+u)\notag\\
 &-\mathbf1_{N\text{ odd}}
      \frac{2N!\zeta(N+1)}{p!}\,\frac{P_p(u)}u\bigg].
 \label{ct:R}
\end{align}
For $N=0$ the first line is the removable expression
$-\pi\tan(\pi u/2)\zeta(1+u)$.

\begin{theorem}[All-index coincident reflected closure]
\label{ct:closure}
The function $\mathcal R_{p,r}$ is holomorphic near $u=0$, and for every
$m,p,r\ge0$,
\begin{equation}\label{ct:moment}
 \boxed{\displaystyle
 \operatorname{FP}\int_0^1\gamma_m^{(p)}(x)K^{(r)}(x)\,dx
       =(-1)^m m!\,[u^m]\mathcal R_{p,r}(u).}
\end{equation}
Thus all these moments have a finite reduction to ordinary Riemann zeta
jets at $N+1$ and generalized harmonic polynomials. In the exceptional
case $N=0$, ordinary Stieltjes constants occur instead.
\end{theorem}
\begin{proof}
Write $g(u,x)=\zeta(1+u,x)-1/u$. The exact local splitting is
\[
 \partial_x^pg(u,x)
   =(-1)^pP_p(u)x^{-p-1-u}+h_p(u,x),\qquad
 h_p(u,x)=\partial_x^pg(u,1+x).
\]
The function $h_p$ is holomorphic jointly near $(u,x)=(0,0)$.
At the upper endpoint the entire factor $\partial_x^pg(u,1-y)$ is
holomorphic jointly near $(0,0)$ in $(u,y)$. The cotangent expansion is
\[
 K^{(r)}(x)=(-1)^rr!x^{-r-1}
 -2\sum_{j\ge1}\zeta(2j)
        \frac{(2j-1)!}{(2j-1-r)!}x^{2j-1-r},
\]
where a term is zero when $2j-1<r$.
In the product of the varying singular power with this expansion,
resonance with $x^{-1-u}$ occurs exactly when
$2j=N+1$, hence exactly when $N$ is odd. Its complete coefficient is
\[
 A_{p,r}(u)
 =-2(-1)^p\frac{N!}{p!}\zeta(N+1)P_p(u).
\]
The meromorphic integral of this local monomial on $(0,b)$ is
$-A_{p,r}(u)b^{-u}/u$. Its coefficientwise coordinate finite-part
generator is $-A_{p,r}(u)(b^{-u}-1)/u$. Their difference is
$-A_{p,r}(u)/u$; therefore the required completion is the meromorphic
integral plus $A_{p,r}(u)/u$.

All other varying-power monomials have denominators nonzero at zero.
Subtract finitely many local Taylor terms at both endpoints so that all
remainders and any fixed finite number of parameter derivatives are
bounded by $Cx^{-\eta}(1+|\log x|^M)$, respectively
$Cy^{-\eta}(1+|\log y|^M)$, with $\eta<1$. These integrals are normally
convergent. The fixed integer endpoint powers from the smooth factors
have analytic coefficients and their coordinate finite parts may be
extracted coefficientwise. This proves holomorphy and identifies all
Taylor coefficients.

The meromorphic integral itself is, by \eqref{ct:J},
\[
 (-1)^pP_p(u)J_r(1+p+u)
 =(-1)^pP_N(u)\pi\cot\frac{\pi(N+1+u)}2\zeta(N+1+u).
\]
The subtraction $-K^{(r)}/u$ when $p=0$ contributes nothing because
$K^{(r)}$ has coordinate finite-part mean zero. Adding the complete
local correction just obtained gives exactly \eqref{ct:R}, and
coefficient extraction in $g$ proves \eqref{ct:moment}.
\end{proof}

Removing only the pole residue would replace $P_p(u)$ by $p!$ and
lose finite terms. In particular, for odd $N$ it would incorrectly
omit the harmonic correction $H_p$ in the next formula.

\subsection{Polygamma specializations and exact finite endpoint anomalies}

Let $H_0=0$. At Stieltjes index zero, \eqref{ct:moment} gives
\begin{equation}\label{ct:zero-index}
 \operatorname{FP}\int_0^1\psi^{(p)}(x)K^{(r)}(x)\,dx
 =\begin{cases}
 \pi^2/2,&N=0,\\
 0,&N>0\text{ even},\\
 2(-1)^{p+1}N!\bigl[\zeta'(N+1)
                +(H_N-H_p)\zeta(N+1)\bigr],&N\text{ odd}.
 \end{cases}
\end{equation}
The first examples are
\begin{align*}
 \operatorname{FP}\int_0^1\psi(x)K'(x)\,dx
   &=-2\zeta'(2)-2\zeta(2),\\
 \operatorname{FP}\int_0^1\psi'(x)K(x)\,dx
   &=2\zeta'(2),\\
 \operatorname{FP}\int_0^1\psi(x)K'''(x)\,dx
   &=-12\zeta'(4)-22\zeta(4),\\
 \operatorname{FP}\int_0^1\psi'(x)K''(x)\,dx
   &=12\zeta'(4)+10\zeta(4),\\
 \operatorname{FP}\int_0^1\psi''(x)K'(x)\,dx
   &=-12\zeta'(4)-4\zeta(4),\\
 \operatorname{FP}\int_0^1\psi'''(x)K(x)\,dx
   &=12\zeta'(4).
\end{align*}
These are consistent with finite-part integration by parts, including
its boundary constants. For example,
\[
 \operatorname{FP}\int_0^1(\psi K)'\,dx=-2\zeta(2).
\]
Indeed the constant terms of $\psi(x)K(x)$ at zero and one are
$3\zeta(2)$ and $\zeta(2)$, respectively. Thus setting the integral
of this ordinary derivative to zero would be incorrect in the present
coordinate convention.

At $p=r=0$, every Stieltjes index has the particularly short formula
\begin{align}
 \operatorname{FP}\int_0^1\gamma_m(x)K(x)\,dx
 ={}&4m!\sum_{k=1}^{\lfloor(m+1)/2\rfloor}
   \frac{(1-2^{-2k})\zeta(2k)\gamma_{m-2k+1}}
                       {(m-2k+1)!}\notag\\
 &-\mathbf1_{2\mid m}\,4m!(1-2^{-m-2})\zeta(m+2).
 \label{ct:all-m-base}
\end{align}
For $m=0,1,2$ this gives
\[
 -\frac{\pi^2}{2},\qquad
 \frac{\pi^2\gamma}{2},\qquad
 \pi^2\gamma_1-\frac{\pi^4}{12}.
\]
The base constant can also be recovered formally from the finite
constant of the separated Stieltjes--cotangent transform in the incoming
triple report. The local proof above establishes that this formal limit
has the stated same-point coordinate normalization, and supplies all
argument-derivative resonances.

For $N>0$ even and $m=1$ there is the simple additional family
\[
 \operatorname{FP}\int_0^1\gamma_1^{(p)}(x)K^{(r)}(x)\,dx
   =\frac{(-1)^pN!\pi^2}{2}\zeta(N+1).
\]
All higher coefficients in \eqref{ct:R} are finite because
$P_j(u)=j!\prod_{a=1}^j(1+u/a)$ is a polynomial.

\subsection{Arbitrarily many coincident reflected factors}

Define derivative polynomials
\[
 D_0(X)=X,\qquad D_{r+1}(X)=-(X^2+\pi^2)D_r'(X).
\]
Then $K^{(r)}(x)=D_r(K(x))$, and $D_r$ has degree $r+1$ and leading
coefficient $(-1)^rr!$. Consequently every polynomial $P(X)$ has a unique
finite decomposition
\begin{equation}\label{ct:poly-decomp}
 P(X)=C+\sum_{r=0}^{\deg P-1}c_rD_r(X).
\end{equation}
The constant is explicitly
\[
 C=\frac{P(i\pi)+P(-i\pi)}2,
\]
since $D_r(\pm i\pi)=0$ for $r\ge1$ and $D_0(\pm i\pi)=\pm i\pi$.
The remaining coefficients follow by descending elimination of the
leading powers, an exact terminating procedure.

\begin{corollary}[Finite coincident reflected algebra]
\label{ct:algebra}
For arbitrary $m,p\ge0$ and any polynomial $P$, the finite part
\[
 \operatorname{FP}\int_0^1\gamma_m^{(p)}(x)P(K(x))\,dx
 =\sum_r c_r(-1)^m m![u^m]\mathcal R_{p,r}(u)
\]
has a finite evaluation in ordinary zeta jets and Stieltjes constants.
The same holds after replacing
$P(K)$ by any finite product of arbitrary derivatives $K^{(r_j)}$.
Equivalently, it holds for any finite product of the reflected factors
\[
 (-1)^{r_j}\psi^{(r_j)}(1-x)-\psi^{(r_j)}(x).
\]
\end{corollary}
\begin{proof}
The reflection formula \cite[Eq.~5.15.6]{DLMFPolyGamma} identifies each balanced polygamma expression
with $K^{(r_j)}$. The derivative-polynomial recurrence converts their
product to a polynomial in $K$. Apply \eqref{ct:poly-decomp} and
Theorem~\ref{ct:closure}. The constant term contributes zero because
$\operatorname{FP}\int_0^1\gamma_m^{(p)}(x)\,dx=0$. For $p=0$ this
follows by integrating the shifted Hurwitz family and completing its
single resonance; for $p\ge1$ it follows directly by taking endpoint
constant terms of $\gamma_m^{(p-1)}$. The smooth endpoint jets coincide
and the singular derivative of $x^{-1}\log^m x$ has no constant term.
\end{proof}

For powers alone, write
$K^q=C_q+\sum_rc_{q,r}K^{(r)}$. A useful recurrence is
\[
 C_0=1,\quad C_1=0,\quad c_{1,0}=1,\qquad
 C_q=-\pi^2 C_{q-2},\qquad
 c_{q,r}=-\frac{c_{q-1,r-1}}{q-1}-\pi^2c_{q-2,r},
\]
with absent coefficients zero. It follows from
$(K^{q-1})'=-(q-1)(K^q+\pi^2K^{q-2})$.
In particular,
\begin{align*}
 K^2&=-K'-\pi^2,\\
 K^3&=\tfrac12K''-\pi^2K,\\
 K^4&=-\tfrac16K'''+\tfrac43\pi^2K'+\pi^4,\\
 K^5&=\tfrac1{24}K^{(4)}-\tfrac56\pi^2K''+\pi^4K.
\end{align*}
Since the coordinate finite-part mean of every $K^{(r)}$ is zero,
\begin{equation}\label{ct:pure-power}
 \operatorname{FP}\int_0^1K(x)^q\,dx
       =\pi^q\cos(q\pi/2).
\end{equation}
The resulting mixed digamma examples include
\begin{align}
 \operatorname{FP}\int_0^1\psi(x)K(x)^2\,dx
     &=2\zeta'(2)+2\zeta(2),\label{ct:cubic-reflected}\\
 \operatorname{FP}\int_0^1\psi(x)K(x)^4\,dx
     &=2\zeta'(4)-\frac{8\pi^2}{3}\zeta'(2)
                       -\frac{109}{3}\zeta(4).\label{ct:quintic-reflected}
\end{align}
There is an all-degree elementary family:
\begin{equation}\label{ct:odd-powers}
 \boxed{\displaystyle
 \operatorname{FP}\int_0^1\psi(x)K(x)^{2h+1}\,dx
       =\frac{(-1)^h\pi^{2h+2}}2\qquad(h\ge0).}
\end{equation}
Besides the polynomial proof, reflection gives a short independent
argument: averaging the integral with its reflected copy yields
$-\tfrac12\operatorname{FP}\int K^{2h+2}$, which is
\eqref{ct:odd-powers} by \eqref{ct:pure-power}.

There is also an exactly reducible family with two unreflected factors:
\begin{equation}\label{ct:two-unreflected}
 \operatorname{FP}\int_0^1\psi(x)^2K(x)^{2h+1}\,dx
   =-\operatorname{FP}\int_0^1\psi(x)K(x)^{2h+2}\,dx.
\end{equation}
To prove it, reflect the left side and use
$\psi(1-x)=\psi(x)+K(x)$. The reflected value is its negative minus
twice the right-hand moment before its displayed minus sign, while
the remaining pure odd power has zero finite-part mean. Thus every
moment in \eqref{ct:two-unreflected} also has a finite ordinary-zeta-jet
reduction by Corollary~\ref{ct:algebra}. At $h=0$ it says
\[
 \operatorname{FP}\int_0^1\psi(x)^2K(x)\,dx
       =-2\zeta'(2)-2\zeta(2).
\]

A direct cubic consequence, retaining the distinction from the open
unreflected cube, is
\begin{equation}\label{ct:cubic-difference}
 \operatorname{FP}\int_0^1
 \bigl[\psi(x)^3-\psi(x)^2\psi(1-x)\bigr]dx
      =2\zeta'(2)+2\zeta(2).
\end{equation}
Indeed expanding $\psi K^2$ gives
$\psi\psi(1-x)^2-2\psi^2\psi(1-x)+\psi^3$, and reflection identifies
the finite parts of the first two mixed cubic monomials. This determines
a difference of cubic moments and leaves their common unreflected
Tornheim coordinate free.

\subsection{Log-Gamma primitives and a harmonic-number family}

The same completed transform supplies the integrated Gamma level.
At $s=0$, the only resonance for $\zeta(s,x)K(x)$ is $x^{-1-s}$.
Thus
\[
 \pi\cot(\pi s/2)\zeta(s)+\frac1s
\]
is the generating germ for the coordinate finite parts of the spectral
jets $\zeta^{[j]}(0,x)K(x)$. Its linear coefficient is
$\zeta''(0)+\zeta(2)/2$, because
$\zeta(0)=-1/2$ and
$\pi\cot(\pi s/2)=2/s-\zeta(2)s+O(s^3)$.
Lerch's identity \cite[Eq.~25.11.18]{DLMFHurwitz} $\zeta'(0,x)=\log\Gamma(x)-\tfrac12\log(2\pi)$
and the zero mean of $K$ give
\begin{equation}\label{ct:loggamma-K}
 \operatorname{FP}\int_0^1\log\Gamma(x)K(x)\,dx
       =\zeta''(0)+\frac12\zeta(2).
\end{equation}
The constant $1/s$ is independent of the parameter, apart from its
pole, so this specialization has no omitted finite derivative of a
resonant amplitude.

For $r\ge1$, coordinate integration by parts gives the particularly
simple primitive law
\begin{equation}\label{ct:loggamma-primitive}
 \operatorname{FP}\int_0^1\log\Gamma(x)K^{(r)}(x)\,dx
       =-\operatorname{FP}\int_0^1\psi(x)K^{(r-1)}(x)\,dx.
\end{equation}
Here the boundary constant difference is zero, and this can be checked
rather than assumed. For $r=1$, the constant of $\log\Gamma(x)K(x)$
at each endpoint is $-\gamma$. For $r\ge2$, use
\[
 \log\Gamma(x)=-\log x-\gamma x+
     \sum_{j\ge2}\frac{(-1)^j\zeta(j)}j x^j,
 \qquad
 \log\Gamma(1-y)=\gamma y+
     \sum_{j\ge2}\frac{\zeta(j)}j y^j.
\]
The constant at both endpoints of
$\log\Gamma(x)K^{(r-1)}(x)$ is
$-(r-1)!\zeta(r)/r$. Logarithmic terms are excluded by the same
constant-term prescription. This proves \eqref{ct:loggamma-primitive}
directly from the cutoff fundamental theorem.
Combining it with \eqref{ct:zero-index} yields
\[
 \operatorname{FP}\int_0^1\log\Gamma(x)K^{(r)}(x)\,dx
 =\begin{cases}
 -\pi^2/2,&r=1,\\
 2(r-1)!\bigl[\zeta'(r)+H_{r-1}\zeta(r)\bigr],&r\ge2\text{ even},\\
 0,&r\ge3\text{ odd}.
 \end{cases}
\]
For any polynomial $P$ with decomposition \eqref{ct:poly-decomp}, its
log-Gamma moment is therefore a finite sum of these evaluations and
\eqref{ct:loggamma-K}, with the constant contribution
$C\log(2\pi)/2$. This is an explicit primitive counterpart of the
coincident reflected algebra.

The even powers have an especially compact expression involving only
ordinary harmonic numbers:
\begin{equation}\label{ct:loggamma-even}
 \boxed{\displaystyle
 \operatorname{FP}\int_0^1\log\Gamma(x)K(x)^{2h}\,dx
 =\frac{(-1)^h\pi^{2h}}2
       \left[\log(2\pi)-H_{2h}+\frac12H_h\right]
       \quad(h\ge0).}
\end{equation}
Here is an independent proof of the normalization. Initially for
$\Re a>2h-1$,
\[
 V_h(a):=\int_0^1\sin(\pi x)^aK(x)^{2h}\,dx
  =\pi^{2h-1}B\left(\frac{a-2h+1}{2},h+\frac12\right).
\]
At $a=0$, all endpoint exponents are even integers; none is $-1$.
Thus this beta continuation is holomorphic there, and its value and
first derivative equal the coordinate finite parts of $K^{2h}$ and
$K^{2h}\log\sin(\pi x)$. The beta identity \cite{DLMFBeta}, Gamma reflection, and the digamma recurrence \cite{DLMFGamma}
give
\[
 V_h(0)=(-1)^h\pi^{2h},\qquad
 V_h'(0)=(-1)^h\pi^{2h}
       \left[-\log2+H_{2h}-\frac12H_h\right].
\]
Finally average the log-Gamma integral with its reflected copy and use
$\log\Gamma(x)+\log\Gamma(1-x)=\log\pi-\log\sin(\pi x)$.
This proves \eqref{ct:loggamma-even} in the same unit coordinates.
For example,
\[
 \operatorname{FP}\int_0^1\log\Gamma(x)K(x)^2\,dx
   =\frac{\pi^2}{2}\bigl[1-\log(2\pi)\bigr],
\]
\[
 \operatorname{FP}\int_0^1\log\Gamma(x)K(x)^4\,dx
   =\frac{\pi^4}{2}\left[\log(2\pi)-\frac43\right].
\]
The beta/reflection proof uses classical formulas; it is included as an
explicit infinite family and as an independent check on the primitive
normalization, without a claim of global priority.

\subsection{Extension to repeated seams at arbitrary shifts}

The same argument completes the separated theorem by rational partial
fractions. Any finite product of $K^{(r_j)}(x-a_j)$ is rational in
$X=e^{2\pi ix}$, bounded at $X=0,\infty$, and has poles only at the
finitely many $e^{2\pi ia_j}$. At each pole, the principal parts of
$K^{(r)}(x-a)$ form a triangular basis. Subtracting them leaves a
constant. Hence a unique finite representation exists as
\[
 C+\sum_a\sum_{r=0}^{M_a-1}c_{a,r}K^{(r)}(x-a),
\]
where $M_a$ is the pole order at that seam. Multiplication by the ordinary
function $\gamma_m^{(p)}(x)$ and coordinate finite-part integration are
linear operations on its local Laurent-log expansion, so this rational
identity remains valid under precisely the stated normalization. The
terms with $a=0$ are evaluated by \eqref{ct:moment}; the other terms use
the already proved separated Stieltjes--cotangent kernel and its
contact-corrected argument derivatives. No limit through colliding
separated distributions is required.

For completeness, if
$I_m(a)=\operatorname{FP}\int_0^1\gamma_m(x)K(x-a)dx$ at $0<a<1$,
$N=p+r$, and
$h_{m,p}=(-1)^m m!e_{m+1}(p)$ for $p\ge1$ (zero for $p=0$),
the existing one-factor contact law \cite{RepoMain,RepoTriple} gives
\[
 \operatorname{FP}\int_0^1\gamma_m^{(p)}(x)K^{(r)}(x-a)dx
       =(-1)^r\bigl[I_m^{(N)}(a)-h_{m,p}K^{(N)}(a)\bigr].
\]
Here $e_j(p)$ is the elementary symmetric polynomial of degree $j$ in
$1,1/2,\ldots,1/p$, and is zero for $j>p$. To check signs, use
$F_{m,p}=D^pG_m+h_{m,p}\delta^{(p)}$ and
$K^{(N)}(-a)=(-1)^{N+1}K^{(N)}(a)$. Thus arbitrary repetitions and
coincidences within the reflected algebra still reduce to finitely many
ordinary Hurwitz-zeta jets, with trigonometric coefficients at the
remaining distinct shifts.

\subsection{Verification and remaining questions}

The accompanying verifier computes independent endpoint Laurent-log
series from the ordinary Hurwitz shift relation and integrates them
term by term on $(0,1/5)$ and $(4/5,1)$. On the central interval it uses
ordinary high-precision quadrature of Stieltjes/polygamma and cotangent
functions. It compares these values against finite coefficients of
\eqref{ct:R}, and separately integrates powers $K^q$ without applying
the derivative-polynomial reduction. Exact symbolic assertions verify
the power reduction recurrence. These are independent floating-point
diagnostics, not interval enclosures; the proof above establishes the
all-index result.

The natural next questions are to classify which products with two
unreflected Stieltjes factors and a reflected polynomial still close
without a Tornheim jet, to derive a conductor-compatible version of the
repeated-seam partial fractions, and to formalize the local coefficient
completion independently of distributional products. The present
results do not imply an arithmetic evaluation of the generic cubic
Tornheim derivative or either fixed-weight Gaussian candidate.

% END sections/03_reflected.tex


% BEGIN sections/04_harmonic.tex
% Manuscript-ready fragment: raw harmonic powers at nonpositive spectral integers.
% Standard amsmath, amssymb, amsthm. Labels/citations have hp: prefix.
\section{Raw harmonic powers: polynomial finite parts and a centered cancellation}
\label{hp:sec:main}

The elementary-symmetric harmonic numerators of the height-one family do
not equal raw powers. For example, its second numerator is
$H_n^2-H_n^{(2)}$, whereas the family studied here contains $H_n^2$ itself.
Let
\begin{equation}\label{hp:def:F}
 H_n=\sum_{j=1}^n\frac1j,\qquad H_0=0,\qquad
 F_p(s,a)=\sum_{n=0}^{\infty}\frac{H_n^p}{(n+a)^s},
 \qquad p\ge1,\quad a>0,\quad \Re s>1.
\end{equation}
The inner harmonic numbers are unshifted; only the outer denominator
moves. All finite parts below mean the constant Laurent coefficient of
this meromorphic continuation in the variable $s$. They are not an
assignment of an ordinary value to the divergent defining series.

Dirichlet series of raw harmonic powers, their poles, and reductions to
Euler sums have already been studied, notably by Tin
\cite{hp:Tin2025}. The polylogarithm and shuffle ingredients are also
classical \cite{hp:IKZ2006}. The contribution developed here is the
continuous outer-shift normalization, its polynomial finite-part law at
every nonpositive integer, the exact cancellation at the centered shift,
and explicit low-weight identities with normalized parameter primitives.
No claim of priority is made for raw-power meromorphic continuation or
for the classical Euler sums used in the reductions.

\subsection{Every principal part has a finite Bernoulli formula}

Write $B_j(a)$ for the Bernoulli polynomial, with $B_1(a)=a-1/2$, and
introduce the formal series
\begin{equation}\label{hp:def:U}
 U(z,a)=-\sum_{j\ge1}\frac{B_j(a)}jz^j.
\end{equation}
Only finitely many coefficients are used; convergence of this formal
asymptotic series is neither required nor asserted.

\begin{theorem}[All poles and all principal coefficients]\label{hp:thm:poles}
The function $F_p(s,a)$ continues meromorphically to the entire $s$-plane,
locally uniformly in $a>0$. Its only possible poles are
$1,0,-1,-2,\ldots$. At $s=1$ its principal part is
\begin{equation}\label{hp:eq:principal-one}
 \sum_{r=0}^p\binom pr\frac{r!\gamma^{p-r}}{(s-1)^{r+1}}.
\end{equation}
At $s=-m$, $m\ge0$, put
\[
 F_p(-m+\varepsilon,a)
 =\sum_{h=1}^p\frac{R_{p,m,h}(a)}{\varepsilon^h}
       +C_{p,m}(a)+O(\varepsilon).
\]
Then every principal coefficient is the finite polynomial
\begin{equation}\label{hp:eq:allresidues}
 \boxed{\quad
 R_{p,m,h}(a)=\frac{p!}{(p-h+1)!}
       [z^{m+1}]\bigl(\gamma+U(z,a)\bigr)^{p-h+1},
 \qquad 1\le h\le p.\quad}
\end{equation}
In particular,
\begin{equation}\label{hp:eq:topresidue}
 R_{p,m,p}(a)=-\frac{p!}{m+1}B_{m+1}(a).
\end{equation}
Every $R_{p,m,h}$ has degree at most $m+1$ in $a$.
\end{theorem}

\begin{proof}
Put $x=n+a$. The classical shifted digamma expansion, obtained for
example from Euler--Maclaurin \cite{hp:DLMF}, is
\begin{equation}\label{hp:eq:digamma-expansion}
 H_n=\log x+\gamma-
       \sum_{j=1}^{J}\frac{B_j(a)}{jx^j}+O(x^{-J-1}).
\end{equation}
It holds uniformly when $a$ ranges over a compact positive interval;
the finite differentiated versions have the same uniformity. Raising a
finite truncation to the fixed power $p$ gives
\[
 H_n^p=\sum_{j=0}^{J}\sum_{r=0}^p
  \binom pr[z^j](\gamma+U(z,a))^{p-r}
                    x^{-j}(\log x)^r
       +O\bigl(x^{-J-1}(1+\log^p x)\bigr).
\]
Subtract these finitely many terms in the defining Dirichlet series.
The remainder is normally convergent in $\Re s>-J$ on compact
subsets. The subtracted sums are the corresponding linear combination
of $(-1)^r\partial_s^r\zeta(s+j,a)$. Arbitrarily large $J$ therefore
provides the meromorphic continuation and justifies every fixed
parameter derivative.

The only principal term of
$(-1)^r\partial_s^r\zeta(s+j,a)$ at $s=1-j$ is
$r!/(s+j-1)^{r+1}$. At $j=0$ this proves
\eqref{hp:eq:principal-one}. At $j=m+1$, set $r=h-1$ to obtain
\eqref{hp:eq:allresidues}. The case $r=p$ has no positive-degree
coefficient in $z$, so the pole order at a nonpositive integer is at
most $p$. Formula \eqref{hp:eq:topresidue} and the degree assertion
follow directly from the Bernoulli-polynomial coefficients.
\end{proof}

\begin{corollary}[Cancellation at every nonpositive even integer]
\label{hp:cor:centered}
For every $p\ge1$, the centered function $F_p(s,1/2)$ is holomorphic
at $s=0,-2,-4,\ldots$. Its only poles are at
$s=1,-1,-3,-5,\ldots$; their orders are exactly $p+1$ at $1$ and
exactly $p$ at each negative odd integer.
\end{corollary}
\begin{proof}
All odd Bernoulli polynomials vanish at $1/2$, including $B_1(1/2)$.
Thus $U(z,1/2)$ contains only even powers of $z$ and every coefficient
in \eqref{hp:eq:allresidues} vanishes when $m$ is even.
At $m=2r-1$, the top coefficient is
$-p!B_{2r}(1/2)/(2r)$, which is nonzero because
$B_{2r}(1/2)=(2^{1-2r}-1)B_{2r}\ne0$. The leading coefficient at
$s=1$ is $p!$ by \eqref{hp:eq:principal-one}.
\end{proof}

The cancellation concerns the entire Laurent principal part, not only
the residue. In the centered variable $x=n+1/2$ it is the exact
consequence of
$H_n=\log x+\gamma+1/(24x^2)-7/(960x^4)+\cdots$ having only even
inverse powers after its logarithm.

\subsection{A finite polylogarithm calculation determines every finite part}

For a positive-index word $\boldsymbol{k}=(k_1,\ldots,k_d)$, use
\[
 \operatorname{Li}_{\boldsymbol{k}}(z)
 =\sum_{n_1>\cdots>n_d\ge1}
       \frac{z^{n_1}}{n_1^{k_1}\cdots n_d^{k_d}}.
\]
Define the finite combination
\begin{equation}\label{hp:def:A}
 A_p(t)=\sum_{\substack{\boldsymbol{k}\text{ a composition of }p}}
     \frac{p!}{k_1!\cdots k_d!}
                  \operatorname{Li}_{\boldsymbol{k}}(e^{-t}).
\end{equation}
Ordering the distinct indices in the product $H_n^p$ proves the exact
identity
\begin{equation}\label{hp:eq:kernel}
 (1-e^{-t})\sum_{n\ge0}H_n^pe^{-nt}=A_p(t),\qquad
 \Gamma(s)F_p(s,a)=\int_0^\infty t^{s-1}K_{p,a}(t)\,dt,
 \quad K_{p,a}(t)=\frac{e^{-at}A_p(t)}{1-e^{-t}}.
\end{equation}
The integral is initially ordinary for $\Re s>1$. The first identity
is a finite stuffle decomposition: an ordered equality block of size
$k_j$ contributes $n_j^{-k_j}$ and its multiplicity is the displayed
multinomial coefficient.

For completeness, the local calculation has an explicit finite
recursion. With $L_{\boldsymbol{k}}(t)=
\operatorname{Li}_{\boldsymbol{k}}(e^{-t})$, set $L_{\varnothing}=1$.
Then
\begin{equation}\label{hp:eq:word-derivatives}
 \frac{d}{dt}L_{k_1,\ldots,k_d}(t)=
 \begin{cases}
  -L_{k_1-1,k_2,\ldots,k_d}(t),& k_1>1,\\
  -L_{k_2,\ldots,k_d}(t)/(e^t-1),& k_1=1.
 \end{cases}
\end{equation}
For a convergent word $k_1\ge2$, the constant is
$L_{\boldsymbol{k}}(0)=\zeta(\boldsymbol{k})$.
For a leading-one word, the integration constant is fixed by the
shuffle relations and $L_1(t)=-\log(1-e^{-t})$, whose constant
coefficient in powers of $\log t$ is zero. This is the usual shuffle
regularization. It is a finite reduction: multiplying a word with one
fewer initial $1$ by $\operatorname{Li}_1$ and shuffling its iterated
integral word isolates the desired initial-one word with a nonzero
integer multiplicity. All other terms have fewer initial $1$'s.
Induction reduces its constant to convergent multiple zeta values.

At a specified Taylor order, \eqref{hp:eq:word-derivatives} uses only
rational operations and finite integrations of logarithmic monomials.
Explicitly, for a polynomial $P$,
\begin{align}
 \int t^{-1}P(\log t)\,dt
   &=\int P(\ell)\,d\ell\big|_{\ell=\log t},\notag\\
 \int t^jP(\log t)\,dt
   &=t^{j+1}\sum_{q=0}^{\deg P}
           \frac{(-1)^qP^{(q)}(\log t)}{(j+1)^{q+1}},
                       \qquad j\ge0.\label{hp:eq:finite-integration}
\end{align}
This supplies an implementable calculation with fixed endpoint
constants, rather than an unspecified asymptotic coefficient.

\begin{theorem}[Polynomial finite parts at all nonpositive integers]
\label{hp:thm:finiteparts}
For fixed $p\ge1$ and $m\ge0$, compute the finite local expansion
\begin{equation}\label{hp:eq:local-kernel}
 K_{p,a}(t)=\sum_{j=-1}^{m}t^j
         \sum_{r=0}^p c_{p,j,r}(a)(\log t)^r
       +O\bigl(t^{m+1}(1+|\log t|^p)\bigr).
\end{equation}
The coefficients are provided by
\eqref{hp:def:A}--\eqref{hp:eq:finite-integration}; they are polynomials
in $a$ of degree at most $j+1$ with coefficients in the algebra of
convergent multiple zeta values. Define
\begin{equation}\label{hp:def:g}
 \frac1{\Gamma(-m+\varepsilon)}
      =\sum_{\ell\ge1}g_{m,\ell}\varepsilon^\ell.
\end{equation}
Then the exact finite part is
\begin{equation}\label{hp:eq:FP-all}
 \boxed{\quad C_{p,m}(a)=
   \sum_{r=0}^p(-1)^rr!\,c_{p,m,r}(a)g_{m,r+1}.\quad}
\end{equation}
In particular $C_{p,m}$ is a polynomial in $a$ of degree at most
$m+1$. Its coefficients lie in the algebra generated by Euler's
constant and convergent multiple zeta values. At the centered regular
points one has the stronger exact identity
\begin{equation}\label{hp:eq:center-value}
 F_p(-2r,1/2)=(2r)!\,c_{p,2r,0}(1/2).
\end{equation}
These centered values lie in the multiple-zeta algebra without requiring
Euler's constant.
\end{theorem}

\begin{proof}
Equations \eqref{hp:eq:word-derivatives} and
\eqref{hp:eq:finite-integration}, with the fixed shuffle constants,
construct the expansion by induction on the weight. The ordinary
convergent expansion of $1/(e^t-1)$ has a simple pole at zero and
radius $2\pi$. Finite integration of its analytic remainders preserves
a remainder bounded by a power of $t$ times a fixed logarithmic
polynomial. Consequently the displayed finite expansion is an actual
local expansion with the stated remainder, locally uniformly in $a$;
there are no negative powers below $t^{-1}$. Multiplication by
$e^{-at}$ proves the degree bound.

Subtract \eqref{hp:eq:local-kernel} inside the integral over $(0,1)$.
The remainder integral, together with the integral over $(1,\infty)$,
is holomorphic in $\Re s>-m-1$. Its fixed spectral derivatives are
obtained by inserting powers of $\log t$, by ordinary domination.
Thus
\begin{equation}\label{hp:eq:Mellin-continuation}
 \Gamma(s)F_p(s,a)=
 \sum_{j=-1}^{m}\sum_{r=0}^p
       \frac{(-1)^rr!c_{p,j,r}(a)}{(s+j)^{r+1}}+J_{p,m}(s,a),
\end{equation}
where $J_{p,m}$ is holomorphic near $s=-m$.
The zero of $1/\Gamma(s)$ at $s=-m$ kills the holomorphic remainder
and all terms with $j\ne m$ when taking the constant Laurent
coefficient. Only the finitely many $j=m$ terms remain, and their
product with \eqref{hp:def:g} is precisely \eqref{hp:eq:FP-all}.

For reference the reciprocal-Gamma coefficients are themselves finite
Bell expressions, obtained from
\begin{align}\label{hp:eq:gamma-generator}
 \frac1{\Gamma(-m+\varepsilon)}
 ={}&(-1)^m m!\,\varepsilon
 \exp\left((\gamma-H_m)\varepsilon\right.\notag\\
 &\left.\hspace{8mm}+
 \sum_{k\ge2}\frac{(-1)^{k+1}\zeta(k)-H_m^{(k)}}k
                                      \varepsilon^k\right).
\end{align}
The recurrence for Gamma proves this formula directly from the Taylor
series of $1/\Gamma(1+\varepsilon)$.
At a centered regular point, the highest possible logarithmic
coefficient in \eqref{hp:eq:Mellin-continuation} would produce a
nonzero pole. Descending triangularly, Corollary~\ref{hp:cor:centered}
therefore forces $c_{p,2r,j}(1/2)=0$ for every $j\ge1$. Since
$g_{2r,1}=(2r)!$, this proves \eqref{hp:eq:center-value}.
\end{proof}

\paragraph{Where the global information returns.}
The theorem evaluates the principal part and constant term. It does
not evaluate all spectral derivatives using the same finite data.
Already the coefficient of $\varepsilon$ receives the generally nonzero
contribution $(-1)^m m!J_{p,m}(-m,a)$ from the holomorphic remainder
in \eqref{hp:eq:Mellin-continuation}, together with the other regular
terms. In particular, differentiating a formula valid only for a
nonpositive integer does not determine the spectral derivative there.
This distinguishes the present finite identities from the higher
harmonic Stieltjes coefficients.

\subsection{The complete expansion through the finite part at zero}

Set
\begin{equation}\label{hp:def:P}
 P_p(L)=p![u^p]\exp(Lu)\Gamma(1+u)e^{\gamma u},\qquad P_0=1.
\end{equation}
Define ordinary convergent Euler sums
\begin{equation}\label{hp:def:E}
 E_{q,k}=\sum_{n\ge1}\frac{H_{n-1}^q}{n^k},
 \qquad q\ge0,\quad k\ge2.
\end{equation}
At $q=0$ the numerator is one. Every such sum is the finite
multiple-zeta combination
\begin{equation}\label{hp:eq:E-MZV}
 E_{q,k}=\sum_{\boldsymbol{r}\vDash q}
       \frac{q!}{r_1!\cdots r_d!}\zeta(k,r_1,\ldots,r_d),
\end{equation}
with the empty composition giving $E_{0,k}=\zeta(k)$.

\begin{theorem}[An exact all-degree zero-resonance identity]
\label{hp:thm:zero}
There are constants $Q_p$ with $Q_1=1/2$ and the finite recurrence
\begin{equation}\label{hp:eq:Q-recurrence}
 Q_p=-pQ_{p-1}+
 \frac12\sum_{k=2}^{p-1}(k-1)\binom p{k+1}E_{p-1-k,k},
                       \qquad p\ge2.
\end{equation}
For every $a>0$ one has
\begin{equation}\label{hp:eq:zero-expansion}
 \boxed{\quad F_p(s,a)=
   \left(\frac12-a\right)\frac{p!e^{\gamma s}}{s^p}
                 +Q_p+O(s).\quad}
\end{equation}
Consequently
\begin{equation}\label{hp:eq:zero-FP}
 C_{p,0}(a)=Q_p+\left(\frac12-a\right)\gamma^p,
                  \qquad F_p(0,1/2)=Q_p.
\end{equation}
Each $Q_p$ has a finite expression using convergent multiple zeta
values of weights at most $p-1$.
\end{theorem}

\begin{proof}
First the leading local polynomial of $A_p$ is $P_p(-\log t)$.
One way to fix it without a regularization guess is to use
\eqref{hp:eq:principal-one} in the Mellin identity. Multiplication by
$\Gamma(1+s)$ turns the elementary pole polynomial at $s=0$ into
exactly the coefficients of
$\Gamma(1+u)e^{\gamma u}$ in \eqref{hp:def:P}.
The word calculation proves existence of the local expansion, so this
identifies its leading polynomial uniquely.
At the centered shift
$e^{-t/2}/(1-e^{-t})=t^{-1}+O(t)$.
Regularity at $s=0$ then forces the coefficient of $t$ in $A_p$ to
have no logarithmic terms, by the same triangular argument as in the
previous theorem. Hence
\begin{equation}\label{hp:eq:A-first}
 A_p(t)=P_p(-\log t)+tQ_p+O(t^2(1+|\log t|^p)).
\end{equation}
The elementary formula for $A_1=-\log(1-e^{-t})$ gives $Q_1=1/2$.

To determine every other $Q_p$, put
\[
 M_{q,k}(t)=\sum_{n\ge1}\frac{H_{n-1}^q}{n^k}e^{-nt}.
\]
The exact increment $H_n^p-H_{n-1}^p$ gives
\[
 A_p=\sum_{k=1}^p\binom pkM_{p-k,k},\qquad
 A_p'=-\frac{pA_{p-1}}{e^t-1}
            -\sum_{k=2}^p\binom pkM_{p-k,k-1}.
\]
For $k\ge2$, $M_{q,k}(t)\to E_{q,k}$ by domination.
Using the preceding identity with $p-1$ also gives the complete
leading logarithmic polynomial
\[
 M_{p-2,1}(t)=\frac{1}{p-1}
 \left(P_{p-1}(-\log t)-
       \sum_{k=2}^{p-1}\binom{p-1}kE_{p-1-k,k}\right)+o(1).
\]
Take the coefficient of $t^0(\log t)^0$ in the derivative identity,
using \eqref{hp:eq:A-first} and
$(e^t-1)^{-1}=t^{-1}-1/2+O(t)$.
The two copies of $P_{p-1}$ cancel as full polynomials. This leaves
\[
 Q_p=-pQ_{p-1}+\sum_{k=2}^{p-1}
 \left(\frac p2\binom{p-1}k-\binom p{k+1}\right)E_{p-1-k,k}.
\]
The bracket is $(k-1)\binom p{k+1}/2$, proving the recurrence.
The local expansions already established justify differentiation and
coefficient comparison here.

It remains to show the particularly simple normalization in
\eqref{hp:eq:zero-expansion}. The coefficient of $t^0$ in the kernel is
$(1/2-a)P_p(-\log t)+Q_p$. If
$d(u)=\Gamma(1+u)e^{\gamma u}$, its $P_p$ contribution to the Mellin
continuation is
\[
 \frac{p!}{\Gamma(s)}[u^p]\frac{d(u)}{s-u}.
\]
Replace $d(u)$ by $d(s)$: the difference quotient is holomorphic in
$(s,u)$ near $(0,0)$ and its multiplication by $1/\Gamma(s)$ is
$O(s)$. The remaining term is
$p!e^{\gamma s}/s^p$. The constant $Q_p$ contributes
$Q_p/(s\Gamma(s))=Q_p+O(s)$. All other local and global terms are
$O(s)$ by the preceding theorem. This proves the result.
\end{proof}

The first five constants are
\begin{align}\label{hp:eq:Q-table}
 Q_1&=\frac12,& Q_2&=-1,\notag\\
 Q_3&=3+\frac12\zeta(2),&
 Q_4&=-12-2\zeta(2)+3\zeta(3),\notag\\
 Q_5&=60+10\zeta(2)-15\zeta(3)+\frac{33}{2}\zeta(4).
\end{align}
For example $E_{1,2}=\zeta(2,1)=\zeta(3)$ and
$E_{2,2}=2\zeta(2,1,1)+\zeta(2,2)=11\zeta(4)/4$.
These classical low-weight reductions follow from shuffle and stuffle
relations. Equation \eqref{hp:eq:zero-expansion} includes all negative
Laurent powers: for example the cubic case is
\[
 F_3(s,a)=\left(\frac12-a\right)
 \left(\frac6{s^3}+\frac{6\gamma}{s^2}+
                    \frac{3\gamma^2}{s}+\gamma^3\right)
         +3+\frac12\zeta(2)+O(s).
\]

\subsection{Explicit formulas at the next two spectral integers}

The finite word recursion gives the following local table. Here
$L=-\log t$, $z_j=\zeta(j)$, and the omitted remainder is
$O(t^4(1+|\log t|^p))$.
\begin{align}\label{hp:eq:A-table}
 A_1(t)={}&L+\frac t2-\frac{t^2}{24}+O(t^4),\notag\\
 A_2(t)={}&L^2+z_2-t-\frac{L}{12}t^2-\frac{t^3}{36}
                         +O(t^4(1+L^2)),\notag\\
 A_3(t)={}&L^3+3z_2L-2z_3+\left(3+\frac{z_2}{2}\right)t\notag\\
 &+\left(-\frac{L^2}{8}-\frac L4-\frac{z_2}{8}-\frac38\right)t^2
           +\frac5{72}t^3+O(t^4(1+|L|^3)),\notag\\
 A_4(t)={}&L^4+6z_2L^2-8z_3L+\frac{27}{2}z_4
             +(-12-2z_2+3z_3)t\notag\\
 &+\left(-\frac{L^3}{6}-\frac{L^2}{2}-\frac{z_2L}{2}-L
                 -\frac{z_2}{2}+\frac{z_3}{3}-\frac34\right)t^2
      +\left(-\frac{z_2}{18}-\frac7{27}\right)t^3
                     +O(t^4(1+|L|^4)).
\end{align}
This table is reproducible with \eqref{hp:eq:word-derivatives}.
For clarity the only weight-four endpoint reductions required are
\[
 \zeta(3,1)=\frac14\zeta(4),\qquad
 \zeta(2,2)=\frac34\zeta(4),\qquad
 \zeta(2,1,1)=\zeta(4).
\]
The leading-one regular constants follow by the indicated shuffle
elimination; for example they are $-5\zeta(4)/4$ for $(1,3)$,
$-3\zeta(4)$ for $(1,2,1)$, and $3\zeta(4)$ for $(1,1,2)$.
These constants, the two differentiation rules, and finite integration
fully determine every entry of \eqref{hp:eq:A-table}.

\begin{corollary}[Uniform square, cube, and fourth-power identities]
\label{hp:cor:minus-one}
For $1\le p\le4$, put $b=a-1/2$ and
\[
 T_1=0,\qquad T_2=-\frac1{12},\qquad
 T_3=\frac18,\qquad T_4=-\frac14.
\]
Then
\begin{equation}\label{hp:eq:minus-one-expansion}
 \boxed{\quad
 F_p(-1+\varepsilon,a)=p!e^{\gamma\varepsilon}
 \left(\frac{1/24-b^2/2}{\varepsilon^p}
               +\frac{b^2}{2\varepsilon^{p-1}}\right)
                  +bQ_p+T_p+O(\varepsilon).\quad}
\end{equation}
In particular,
\begin{equation}\label{hp:eq:minus-one-FP}
 C_{p,1}(a)=\frac{\gamma^p}{24}+T_p+bQ_p
                 +\frac{b^2}{2}(p\gamma^{p-1}-\gamma^p).
\end{equation}
The centered regular values at $s=-2$ are
\begin{align}\label{hp:eq:minus-two-table}
 F_1(-2,1/2)&=-\frac1{24},&
 F_2(-2,1/2)&=\frac1{36},\notag\\
 F_3(-2,1/2)&=-\frac19-\frac{\zeta(2)}{24},&
 F_4(-2,1/2)&=\frac{13}{27}+\frac{\zeta(2)}{18}
                                  -\frac{\zeta(3)}4.
\end{align}
\end{corollary}

\begin{proof}
Multiply \eqref{hp:eq:A-table} by
$e^{-at}/(1-e^{-t})$ and apply \eqref{hp:eq:FP-all} at $m=1$.
Together with the Bernoulli formula for the principal part this gives
\eqref{hp:eq:minus-one-expansion}. Equivalently, the centered calculation
uses $e^{-t/2}/(1-e^{-t})=t^{-1}-t/24+O(t^3)$ and the
reciprocal-Gamma expansion at $-1$; the exact parameter transport
proved next supplies its entire $b$-dependence.
For $m=2$, no logarithmic term survives at the centered shift, and
\eqref{hp:eq:center-value} says to multiply the constant coefficient
of $t^2$ by $2$. This coefficient is
$[t^3]A_p-Q_p/24$. Inserting the displayed table gives all four
values in \eqref{hp:eq:minus-two-table}.
\end{proof}

\subsection{Exact shift derivatives and normalized antiderivatives}

Let $R_{p,m}(a)=R_{p,m,1}(a)$ be the simple-pole coefficient in
\eqref{hp:eq:allresidues}. It is given without a polylogarithm
calculation by
\begin{equation}\label{hp:eq:residue-poly}
 R_{p,m}(a)=[z^{m+1}](\gamma+U(z,a))^p.
\end{equation}

\begin{theorem}[Parameter transport and finite primitives]
\label{hp:thm:transport}
For $m\ge1$,
\begin{equation}\label{hp:eq:transport}
 C_{p,m}'(a)=mC_{p,m-1}(a)-R_{p,m-1}(a).
\end{equation}
At $m=0$ the corresponding identity is $C_{p,0}'(a)=-\gamma^p$.
For every $m\ge0$ and positive $a,a_0$ one consequently has the
normalized antiderivative identity
\begin{equation}\label{hp:eq:primitive}
 \boxed{\quad
 \int_{a_0}^a C_{p,m}(x)\,dx=
 \frac{C_{p,m+1}(a)-C_{p,m+1}(a_0)
              +\int_{a_0}^aR_{p,m}(x)\,dx}{m+1}.\quad}
\end{equation}
The last integral is a finite polynomial integral explicitly specified
by \eqref{hp:eq:residue-poly}; it introduces no new special function
or undetermined normalization.
\end{theorem}
\begin{proof}
Normal convergence in the defining half-plane gives
$\partial_aF_p(s,a)=-sF_p(s+1,a)$, and the normally convergent
continuations extend this identity meromorphically. Take the constant
coefficient at $s=-m+\varepsilon$.
The factor $m-\varepsilon$ contributes
$mC_{p,m-1}-R_{p,m-1}$, proving \eqref{hp:eq:transport}.
At $m=0$, the residue at $s=1$ is $\gamma^p$ by
\eqref{hp:eq:principal-one}. Integrating the relation with $m+1$
gives \eqref{hp:eq:primitive} and fixes its additive constant.
\end{proof}

For example, substituting \eqref{hp:eq:zero-FP} and
$R_{p,0}(a)=p(1/2-a)\gamma^{p-1}$ into
\eqref{hp:eq:transport} proves the quadratic shift dependence
\eqref{hp:eq:minus-one-FP} from its centered value. This provides an
independent exact check on the $a$-polynomial obtained by Mellin
coefficient extraction.

\paragraph{Extension and remaining questions.}
A polynomial in finitely many generalized harmonic numbers
$H_n^{(r)}$ of positive integer orders $r$ admits the same finite
ordered-equality-block expansion.
Replacing \eqref{hp:def:A} by its finite multiple-polylogarithm
combination (including a constant when the polynomial has one)
therefore preserves the finite-part theorem, polynomial
shift dependence, and normalized parameter transport. Its pole orders
must be read from its own finite asymptotic expansion; the particular
centered cancellation for raw powers does not automatically extend to
arbitrary generalized-harmonic products.
Three concrete questions remain: determine how far the $Q_p$ reduce
within specified multiple-zeta coordinate spaces at higher weight;
produce compact all-degree formulas for the centered values
$F_p(-2r,1/2)$ beyond the word recursion; and identify combinations of
first regular spectral coefficients whose global Mellin remainders
cancel. The last question is distinct from the finite-part identities
proved here.

% END sections/04_harmonic.tex


% BEGIN sections/05_cayley.tex
% Integration-ready contribution. Requires amsmath, amssymb, amsthm.
\providecommand{\shuffle}{\mathbin{\amalg}}
\section{The complete Cayley ideal does not close the bounded weight-seven search}
\label{gf:sec}

The pinned manuscript leaves the mixed Gaussian identities for $S_6$ and
$S_8$ conjectural. Two recent source contributions sharpened the algebraic
problem in complementary ways. One supplied an exact integer functional
separating the $S_6$ target from a specified collection of $5131$ relations
of depth at most four \cite{gf:Bounded}. The other constructed a
depth-preserving retraction for the \emph{full} shuffle ideal generated by
convergent Cayley identities \cite{gf:Depth}. It explicitly proposed
projecting the bounded rows by that retraction as a next proof attempt.

We complete that proposed calculation. The inherited integer functional
also annihilates the intersection of the \emph{entire} Cayley shuffle ideal
with the weight-seven space of depth at most four. This includes all
Cayley rows and all their shuffle multiples, even when their construction
uses higher-depth intermediate words. Thus completing the Cayley side
alone cannot prove $S_6$ from the stated bounded additional relations.
The numerical $S_6$ identity remains open. This is a result about an
explicit formal presentation; it is not nonvanishing of its evaluated
period residual, and it gives no period-independence assertion.

\subsection{The target and the exact formal presentation}

Use the convention
\[
 \operatorname{Li}_{r,s}(u,v)=\sum_{n>m\geq1}\frac{u^n v^m}{n^r m^s},
 \qquad g_{r,s}=\operatorname{Im}\operatorname{Li}_{r,s}(i,1),
 \qquad S_p=\sum_{n\geq0}\frac{(-1)^nH_n}{(2n+1)^p}.
\]
Let $L=\log2$, and write the conjectural residual in its already reconciled
short coordinate form:
\begin{align}
 R_6={}&S_6+\frac{12334}{527}g_{6,1}+\frac{384}{31}g_{5,2}
       +\frac{2136}{527}g_{4,3}-\frac{3179\pi^7}{5713920}\notag\\
      &+\frac{801}{2108}\beta(4)\zeta(3)+2\beta(6)L.
 \label{gf:eq:residual}
\end{align}
The conjecture is $R_6=0$.

For formal words take the positive forms
\[
 z=\frac{dt}{t},\quad c=\frac{dt}{1-t},\quad
 a=\frac{i\,dt}{1-it},\quad b=-\frac{dt}{1+t},\quad
 \bar a=-\frac{i\,dt}{1+it}.
\]
Their iterated integrals $H(w)$ are written outer letter first on $[0,1]$.
An admissible nonempty word begins with a letter other than $c$ and ends
with a letter other than $z$. Let $\mathcal A$ be the graded rational
shuffle algebra on the admissible words and the empty word. Conjugation
$K$ interchanges $a,\bar a$ and fixes $z,c,b$. The superscript $-$ denotes
its negative eigenspace, and $F_D\mathcal A$ is the span of words with
at most $D$ nonzero letters. This is the usual series-depth filtration.

The map induced by $t\mapsto(1-t)/(1+t)$ is
\[
 C(v_1\cdots v_n)=\tau(v_n)\cdots\tau(v_1),\qquad
 \begin{array}{c|ccccc}
 v&z&c&a&b&\bar a\\ \hline
 \tau v&c-b&z+b&b-\bar a&b&b-a.
 \end{array}
\]
Here $\tau=-\phi^*$ incorporates the reversal of orientation. Every
transformed admissible word is admissible, and change of variables gives
$H(Cw)=H(w)$. Consequently the period map annihilates the graded,
$K$-stable shuffle ideal
\begin{equation}
 \mathcal J=\langle f-Cf:f\in\mathcal A\rangle_{\shuffle}.
 \label{gf:eq:ideal}
\end{equation}
No restriction on weight, depth, or multiplier is made in this definition.

In the following polynomial, juxtaposition is concatenation and $\star$
is the usual series stuffle product, transferred to words by cumulative
colors. Define
\begin{align}
 P_6={}&128588z^5a\bar a+3138084z^5aa
          +1592832z^4aza+521184z^3az^2a-216172z^6a\notag\\
       &+48861\bigl((z^3a)\star(z^2c)\bigr)
          -257176\bigl((z^5a)\star b\bigr),\qquad
 T=\frac{P_6-KP_6}{2}.
 \label{gf:eq:target}
\end{align}
This is the exact target used by the earlier bounded certificate.
It belongs to $F_2\mathcal A_7^-$. The elementary harmonic split gives
$S_6=g_{6,1}+\operatorname{Im}\operatorname{Li}_{6,1}(i,-1)$, while
$\beta(7)=61\pi^7/184320$. Using also $H(b)=-L$ gives
\begin{equation}
 H(T)=128588\,iR_6.
 \label{gf:eq:target-value}
\end{equation}
Both products in \eqref{gf:eq:target} use convergent inputs, so their
stuffle expansion requires no regularization.

For precision, $\mathscr R_4\subset F_4\mathcal A_7^-$ denotes the
\emph{same} inherited finite row family as in \cite{gf:Bounded}. Its
construction is retained in the delivered independent enumerator:
\begin{enumerate}
\item All weight-seven shuffle-minus-stuffle products with sum of input
      depths at most four. Inputs are admissible except for the permitted
      single factor $\operatorname{Li}_1(1)$; the standard one-divergence
      comparison cancels its divergent term.
\item All convergent level-two-to-level-four distribution rows and their
      products with admissible complementary-weight words, provided the
      sum of depths is at most four. Products are expanded by stuffle.
\item The already proved imaginary $S_2$ and $S_4$ rows multiplied by real
      parts of all admissible complementary-weight words of depth at most
      two, again expanded by stuffle.
\item The inherited balanced Cayley rows and their specified products.
      In the adapted alphabet below, the seed $u$ begins with a letter
      other than $x$, ends with a letter other than $z$, and satisfies
      $\min(n_z(u),n_x(u))\geq1$. A multiplier $v$ is allowed when
      $|u|-\min(n_z(u),n_x(u))+\operatorname{dep}(v)\leq4$ and
      $|u|+|v|=7$. Its product is expanded by stuffle.
\end{enumerate}
Imaginary projection, rational normalization, and duplicate removal give,
in this order, $3460$, $1394$, $40$, and $237$ new nonzero rows, respectively:
$5131$ rows in total. These counts specify a formal enumeration, not
dimensions of periods. The fourth family may be retained even though the
new assertion explicitly imposes the complete ideal \eqref{gf:eq:ideal}.

\begin{theorem}[Full-Cayley extension of the weight-seven obstruction]
\label{gf:thm:main}
With the formal target and row family just specified,
\begin{equation}
 \boxed{\qquad
 T\notin\mathcal J_7^-+\operatorname{span}_{\mathbb Q}\mathscr R_4.
 \qquad}
 \label{gf:eq:nonmembership}
\end{equation}
In particular the inherited finite proof search remains unsuccessful even
after \emph{all} convergent Cayley relations and arbitrary shuffle products
of them have been included.
\end{theorem}

The distinction between the full ideal and its bounded intersection is
essential. We verify a finite basis in $F_4\mathcal A_7^-$, but the theorem
allows $\mathcal J$ in the entire admissible word algebra. The next two
subsections justify why those finite checks imply that stronger assertion.

\subsection{Why the depth-preserving projector controls the entire ideal}

We state and justify the precise prerequisite from \cite{gf:Depth}, so the
logic of the extension does not depend on a finite-rank extrapolation.
Introduce
\[
 x=c-b,\qquad y=a-\bar a,\qquad h=a+\bar a-b,
 \qquad \mathcal E=\{z,y,h,b,x\}.
\]
Then
\[
 c=x+b,\qquad a=(y+h+b)/2,\qquad
 \bar a=(-y+h+b)/2.
\]
In this basis admissibility means first letter different from $x$ and last
letter different from $z$. Depth is $|w|-n_z(w)$, conjugation acts by
$(-1)^{n_y(w)}$, and $C$ reverses a word, interchanges $z,x$, and
multiplies by $(-1)^{n_h(w)}$.

Let $\widetilde{\mathcal A}$ be the full shuffle algebra, including
nonadmissible formal words. It is used here only algebraically. With the
order $z<y<h<b<x$, every Lyndon word of length at least two is admissible.
The shuffle-polynomial theorem \cite{gf:Radford} consequently gives
\[
 \widetilde{\mathcal A}=\mathcal A[z,x].
\]
Let $R_0$ be the algebra retraction setting the two \emph{shuffle polynomial}
variables $z,x$ to zero. It preserves both letter counts, commutes with
$C,K$, and is explicitly
\begin{equation}
 R_0(w)=\sum_{u=0}^{l}\sum_{v=0}^{r}(-1)^{u+v}
          x^u\shuffle w[u:|w|-v]\shuffle z^v,
 \label{gf:eq:R0}
\end{equation}
where $l,r$ are the initial $x$ run and terminal $z$ run. Deleting an
initial $x$ and deleting a terminal $z$ are commuting shuffle derivations;
their finite Taylor substitutions prove \eqref{gf:eq:R0}. The factorials
cancel because shuffle powers of a single letter are factorial multiples
of its concatenation powers. In particular this is not literal erasure
of letters inside a word, and assigns no analytic value to divergent
integrals.

For nonempty $w$ put
\begin{equation}
 e(w)=\sum_{k=1}^{|w|}\frac{(-1)^{k-1}}k
       \sum_{w=w_1\cdots w_k}w_1\shuffle\cdots\shuffle w_k,
 \qquad E=R_0e,
 \label{gf:eq:E}
\end{equation}
where the blocks are consecutive and nonempty. The map $e$ is the first
Eulerian logarithm. It annihilates products of positive-degree factors and
equals the identity modulo such products \cite{gf:Patras}. These facts
also follow from the convolution logarithm of the identity character
into the commutative shuffle algebra: its logarithm is an infinitesimal
character, and the term $k=1$ gives the congruence. Thus $e^2=e$.
Reversal permutes consecutive cuts, so $e$ commutes with $C$; it also
commutes with $K$ and preserves both counts.

The subspace $U=E(\mathcal A)$ maps isomorphically onto the
indecomposables of $\mathcal A$, and multiplication identifies
$\mathcal A$ with $\operatorname{Sym}(U)$. To see this directly, $E$ kills
positive-degree products and is the identity modulo them. Any homogeneous
lift of an indecomposable basis therefore gives polynomial generators,
by a triangular change from the Lyndon generators. The lifts are bigraded
by $(n_z,n_x)$, and $C$ exchanges these two degrees.

For a homogeneous $u\in U$ of degree $(p,q)$ choose the orientation
\[
 Q_{\uparrow}u=
 \begin{cases}
 u,&p>q,\\
 Cu,&p<q,\\
 (u+Cu)/2,&p=q.
 \end{cases}
\]
On each pair of different bidegrees this identifies the two transported
generators. On the diagonal it retains the positive eigenspace of $C$ and
kills its negative eigenspace. Extending multiplicatively gives an
algebra retraction whose kernel is exactly $\mathcal J$: on the quotient
all $u-Cu$ vanish, hence $C$ acts identically on every polynomial, and
conversely $u-Cu$ are among the ideal generators. With
$E_{\uparrow}=Q_{\uparrow}E$, the retraction is the finite formula
\begin{equation}
 \Pi_{\uparrow}(w)=\sum_{k=1}^{|w|}\frac1{k!}
       \sum_{w=w_1\cdots w_k}
       E_{\uparrow}(w_1)\shuffle\cdots\shuffle E_{\uparrow}(w_k),
 \qquad \Pi_{\uparrow}(1)=1.
 \label{gf:eq:Pi}
\end{equation}
Convolution exponentiation of \eqref{gf:eq:E}, followed by $R_0$ and
the multiplicative orientation, proves this expression. For completeness,
$E(w)\in U$ even if a block $w$ is not admissible: in degree at least two,
the polynomial decomposition $\mathcal A[z,x]$ makes $w-R_0w$ a sum of
positive-degree products, killed by $e$; in degree one, $R_0e$ kills the
two extra generators directly.

The orientation never decreases the number of $z$'s in any nonzero
output. Shuffle multiplication adds their counts. Consequently, on
$\mathcal A$,
\begin{equation}
 \Pi_{\uparrow}^2=\Pi_{\uparrow},\quad
 \ker\Pi_{\uparrow}=\mathcal J,\quad
 \Pi_{\uparrow}K=K\Pi_{\uparrow},\quad
 \Pi_{\uparrow}(F_D\mathcal A)\subseteq F_D\mathcal A.
 \label{gf:eq:projector-properties}
\end{equation}
These are all-weight algebraic statements. Notice that
$C\Pi_{\uparrow}=\Pi_{\uparrow}$ is neither required nor generally true.

The verifier evaluates \eqref{gf:eq:E} using a permutation formula; its
equivalence with consecutive cuts is also elementary in every weight.
For a permutation of the $n$ positions let $d$ be the number of descents
of its inverse. A shuffle of consecutive blocks produces that permutation
exactly when cuts are placed at those $d$ required positions, with any
additional cuts chosen freely. Its total coefficient in \eqref{gf:eq:E}
is therefore
\[
 \sum_{k=d+1}^n\frac{(-1)^{k-1}}k
      \binom{n-1-d}{k-1-d}
   =(-1)^d\int_0^1 t^d(1-t)^{n-1-d}\,dt
   =\frac{(-1)^d}{n\binom{n-1}{d}}.
\]
This is precisely the implemented coefficient. Repeated letters simply
combine the integer multiplicities of equal output words. Thus the
finite audits of the two formulas are convention checks; their equality
at weight seven does not depend on extrapolating those audits.

\begin{lemma}[The full ideal intersection has a finite spanning certificate]
\label{gf:lem:intersection}
If $\mathcal B_4$ is any rational basis of $F_4\mathcal A_7^-$, then
\[
 \mathcal J_7^-\cap F_4\mathcal A_7
   =\operatorname{span}_{\mathbb Q}
       \{w-\Pi_{\uparrow}(w):w\in\mathcal B_4\}.
\]
\end{lemma}
\begin{proof}
Every displayed difference belongs to the full ideal by the kernel
property and belongs to the bounded space by depth preservation.
Conversely, for $f$ in the intersection, $\Pi_{\uparrow}f=0$, so
$f=(1-\Pi_{\uparrow})f$. Expanding $f$ in the finite basis proves the
opposite inclusion. No restriction on the construction of $f$ inside
$\mathcal J$ has been imposed.
\end{proof}

\subsection{The integer certificate and the proof}

Use the adapted-letter encoding $Z=0,Y=1,H=2,B=3,X=4$. Enumerate all
length-seven words lexicographically, retaining exactly those with first
letter different from $X$, last letter different from $Z$, at least three
$Z$'s, and an odd number of $Y$'s. This gives a basis $\mathcal B_4$ of
$F_4\mathcal A_7^-$ with $2546$ elements. The exact-depth counts are
$1,35,390,2120$. They can be checked directly from the enumeration, or
by choosing the nonzero positions and projecting to odd $Y$ parity.

The certificate's \texttt{separator.json} supplies an integer functional
$\lambda$ on this basis. It is obtained by expressing the inherited
functional of \cite{gf:Bounded} in adapted coordinates; the functional
itself is not being claimed as a newly discovered vector. The new fact
is the second row of the following exact certificate:
\begin{align}
 \lambda(R)&=0 &&(R\in\mathscr R_4),\label{gf:eq:rows-zero}\\
 \lambda(w-\Pi_{\uparrow}w)&=0 &&(w\in\mathcal B_4),
       \label{gf:eq:full-ideal-zero}\\
 \lambda(T)&=-\mathcal N\ne0,\label{gf:eq:pairing}
\end{align}
where the positive integer is
\[
 \begin{aligned}
 \mathcal N={}&186660627289236812620583691130496682\cdot10^{36}\\
              &+78843750753489297279472069700354208.
 \end{aligned}
\]
The integer is an exact functional value, not a numerical residual of the
polylogarithm conjecture. Its decimal expansion is also recorded in the
certificate JSON.

\begin{proof}[Proof of Theorem~\ref{gf:thm:main}]
The standalone verifier regenerates the complete inherited row family,
checks \eqref{gf:eq:rows-zero}, and independently reconstructs each
$\Pi_{\uparrow}(w)$ by \eqref{gf:eq:R0}--\eqref{gf:eq:Pi}. It verifies
\eqref{gf:eq:full-ideal-zero} for all $2546$ basis words, and reconstructs
\eqref{gf:eq:target} before checking \eqref{gf:eq:pairing}. All arithmetic
is integral or rational. Thus Lemma~\ref{gf:lem:intersection} shows that
$\lambda$ annihilates $\mathcal J_7^-\cap F_4\mathcal A_7$.

If, contrary to the theorem, $T=j+r$ with $j\in\mathcal J_7^-$ and
$r\in\operatorname{span}\mathscr R_4$, then $T,r\in F_4\mathcal A_7^-$
implies $j=T-r\in F_4\mathcal A_7^-$. Consequently $\lambda(j)=0$;
also $\lambda(r)=0$ by \eqref{gf:eq:rows-zero}. This contradicts
$\lambda(T)\ne0$ and proves \eqref{gf:eq:nonmembership}.
\end{proof}

For reproducibility, execute
\begin{quote}\ttfamily
python certificates/cayley/code/verify\_complete\_cayley.py
\end{quote}
from the package root. It requires only Python's
standard library. The verifier checks SHA-256 digests of the two retained
source modules and the functional data before use. It regenerates $5131$
rows, verifies $2546$ complete-ideal basis differences, and checks the
nonzero target pairing. There are $2300$ nonzero projector differences;
no claim of their linear independence is made. Additional finite audits
compare the Eulerian permutation formula against consecutive cuts on all
$780$ words through weight four and check $2150$ shuffle products, as
well as retraction, conjugation, and depth conventions. These audits
support the implementation; the preceding algebraic proof supplies the
all-weight projector properties and the passage to the full ideal.

\subsection{Consequences for further research}

The immediate computational recommendation is sharper than in the
inspected sources. Completing Cayley symmetry while retaining precisely
$\mathscr R_4$ does not settle $S_6$. A successful continuation must change
the additional relation input: for example by admitting suitable
higher-depth double-shuffle or distribution rows, by adding another
functional equation, or by deriving a uniform generating identity that
implies the target. A larger numerical-precision calculation does not
address this missing algebraic information.

The theorem does not prove that every conceivable derivation must pass
through depth five. It excludes exactly the full ideal plus the specified
finite additional span; another depth-four family could carry information
outside that span. Likewise, no conclusion for the revised $S_8$ candidate
follows merely by analogy. Useful next questions are to identify the
lowest-weight additional row whose projected class pairs nontrivially
with $\lambda$, to test whether selected depth-five double-shuffle rows
suffice, and to find a conceptual functional equation explaining the
whole even-index family. These are exact proof-search questions, while
$R_6=0$ itself remains the separate numerical-period conjecture.

% END sections/05_cayley.tex


% BEGIN sections/06_audit.tex
\section{Corrections, source audit, and integration safeguards}
\label{audit:section}

\subsection{What was inspected and what was not established}

The source inspection covered the canonical manuscript's status and
relevant harmonic, finite-part, Stieltjes, Gamma, and Gaussian chapters;
the incoming archives were extracted and their relevant theorems,
normalizations, future questions, and verification notes compared.
The most direct dependencies are the ordered-Hurwitz,
triple-Stieltjes--Tornheim, coincident-Stieltjes, directional-zeta, and
resonance-coordinate reports. Two additional pinned formal-algebra
sources supply the finite separator and the depth-preserving Cayley
projector. This is a focused mathematical audit of the material used
here, not a line-by-line certification of the entire multi-volume
manuscript or every incoming claim.

No false theorem was identified in the selected incoming cubic and
coincident claims. In particular, the two previously supplied forms of
the diagonal cubic Tornheim relation are compatible after the zeta
functional equation is applied. Their common third-derivative coordinate
does not disappear merely because a reflected cubic difference is
evaluated in Section~\ref{ct:section}. The older corrections to the
rejected $S_8$ coefficient vector and related proof-status warnings are
already reflected in the inspected manuscript; they should not be
presented as newly discovered errors.

There is one definite external formula correction, proved in the next
subsection. The remaining items below are precise safeguards for
integrating the present results. They are not allegations that the
inspected source texts already make each prohibited inference.


% BEGIN sections/06b_tin_sign.tex
% Optional audit subsection; all labels/citations have hp: prefix.
\subsection{A sign correction and its full logarithmic extension}
\label{hp:sec:tin-audit}

For this comparison use the unshifted raw-power harmonic zeta function
\[
 \mathcal H_p(s)=\sum_{n\ge1}\frac{H_n^p}{n^s},\qquad p\ge0,
 \quad \mathcal H_0(s)=\zeta(s),
\]
and define its regular Laurent coefficients by
\begin{equation}\label{hp:eq:tin-convention}
 \mathcal H_p(1+u)=\operatorname{PP}_{u=0}\mathcal H_p(1+u)
       +\sum_{j\ge0}\frac{(-1)^j}{j!}
                         \widetilde\gamma_{p,j}u^j.
\end{equation}
Thus $\widetilde\gamma_{0,j}=\gamma_j$, the ordinary Stieltjes
constant, and $\widetilde\gamma_{p,j}$ corresponds to Tin's
${}_p\widetilde\gamma_j$ in equation (14) of
\cite{hp:Tin2025}.

\begin{proposition}[Convergent identities for every logarithmic moment]
\label{hp:prop:tin-all-log}
For all integers $p,j\ge0$,
\begin{equation}\label{hp:eq:tin-all-log}
 \boxed{\quad
 \sum_{n\ge1}\frac{H_n^p(H_n-\log n-\gamma)\log^j n}{n}
   =\widetilde\gamma_{p+1,j}
        -\widetilde\gamma_{p,j+1}
        -\gamma\widetilde\gamma_{p,j}.\quad}
\end{equation}
The series is an ordinary absolutely convergent series.
\end{proposition}
\begin{proof}
Since $H_n-\log n-\gamma=1/(2n)+O(n^{-2})$, the Dirichlet
series with numerator $H_n^p(H_n-\log n-\gamma)$ and every fixed
logarithmic derivative converge normally in $\Re s>0$. In $\Re s>1$
its sum equals
\[
 \mathcal H_{p+1}(s)+\mathcal H_p'(s)-\gamma\mathcal H_p(s),
\]
so every principal part cancels at $s=1$. Comparing the coefficient
of $u^j$ in \eqref{hp:eq:tin-convention}, and then multiplying by
$(-1)^j j!$, proves \eqref{hp:eq:tin-all-log}.
\end{proof}

With the convention \eqref{hp:eq:tin-convention}, the first term on
the right of equation (22) of version 1 of \cite{hp:Tin2025},
printed page 5, must be
$-{}_p\widetilde\gamma_1$, in place of the printed
$+{}_p\widetilde\gamma_1$. This corrects that sign; no other term
of that equation is being audited here. The correction agrees with
the paper's own $p=1$ equation (11). Indeed,
\[
 \widetilde\gamma_{1,0}=\frac{\gamma^2+\zeta(2)}2,
 \qquad
 \widetilde\gamma_{2,0}=\frac{\gamma^3+5\zeta(3)}3
\]
give the two illustrative cases
\begin{align}
 \sum_{n\ge1}\frac{H_n-\log n-\gamma}{n}
   &=\frac{\zeta(2)-\gamma^2}{2}-\gamma_1,
       \label{hp:eq:tin-pzero}\\
 \sum_{n\ge1}\frac{H_n(H_n-\log n-\gamma)}{n}
   &=\frac53\zeta(3)-\frac\gamma2\zeta(2)
       -\frac{\gamma^3}{6}-\widetilde\gamma_{1,1}.
       \label{hp:eq:tin-pone}
\end{align}
The two displayed regular constants follow directly from
$\sum_{n\le N}H_n/n=(H_N^2+H_N^{(2)})/2$ and
$\sum_{n\le N}H_n^2/n=H_N^3/3+
\sum_{n\le N}H_n/n^2-H_N^{(3)}/3$, together with
$\sum_{n\ge1}H_n/n^2=2\zeta(3)$.

% END sections/06b_tin_sign.tex


\subsection{Five safeguards with mathematical consequences}

\paragraph{Retain the full resonant amplitude.}
In \eqref{ct:R}, for odd $N=p+r$, the counterterm contains the entire
polynomial $P_p(u)/u$. Replacing it by its principal pole alone changes
the constant and higher coefficients. The difference is observable
already in
\[
 \FP\int_0^1\psi'(x)K(x)\,dx=2\zeta'(2),\qquad
 \FP\int_0^1\psi(x)K'(x)\,dx=-2\zeta'(2)-2\zeta(2).
\]
The harmonic correction in \eqref{ct:zero-index} is part of the identity,
not a removable convention. A source implementation should store the
completed germ, rather than store its pole residue and reconstruct a
finite part heuristically.

\paragraph{A total derivative can have a nonzero finite-part integral.}
With the fixed endpoint coordinates,
\[
 \FP\int_0^1 f'(x)\,dx=\CT_{x\to1}f(x)-\CT_{x\to0}f(x)
\]
whenever the endpoint Laurent-log expansions exist. Thus
$\FP\int_0^1(\psi K)'=-2\zeta(2)$, although the Gamma primitive
in \eqref{ct:loggamma-primitive} has zero boundary constant difference.
These two cases must be checked separately. A blanket rule that all
regularized total derivatives integrate to zero would be false.

\paragraph{Coordinate cancellation is not arithmetic minimality.}
The locus $b=c=d$ in Section~\ref{oq:section} is exactly the locus
where the coefficients of the named additional coordinates in the
displayed normal form vanish. The three-ray reconstruction is an
invertible coordinate calculation. Neither statement proves that
$\eta,\rho,\kappa,T$ are arithmetically independent over a specified
field of ordinary constants. Additional functional relations could
alter a later minimal presentation. The proposed status language is
``proved spanning normal form with an exact cancellation locus in that
form.''

\paragraph{A finite part and a spectral derivative use different data.}
At $s=-m$, multiplication by $1/\Gamma(s)$ kills the holomorphic
Mellin remainder in the constant Laurent coefficient of the raw-power
series. It generally does not kill its contribution to the coefficient
of $s+m$. Consequently, the polynomial formula for $C_{p,m}(a)$
cannot be differentiated with respect to the discrete label $m$ to
obtain a spectral derivative. Actual argument differentiation is valid
and gives \eqref{hp:eq:transport}, including its residue term. This
distinction is essential for any primitive or higher-Stieltjes
extension.

\paragraph{The complete Cayley result still has a fixed additional family.}
The ideal $\mathcal J$ in Section~\ref{gf:sec} is unrestricted, but
$\mathscr R_4$ is the enumerated additional family. The new theorem
rules out every construction within
$\mathcal J_7^-+\operatorname{span}_{\mathbb Q}\mathscr R_4$.
It does not rule out additional double-shuffle, distribution, or other
functional information. Nor does it prove that the actual period
residual is nonzero or that every possible proof must pass through
depth five. The integer pairing is a formal algebraic certificate,
not a lower bound for a period.

\subsection{Proposed canonical status updates}

The following replacements are suitable for the research-status record.
\begin{enumerate}
\item Mark the depth-four independent-slope request as answered by a
spanning formula with convergent coordinates and a formal cancellation
locus. Leave arithmetic minimality and a classification modulo all
functional relations open.
\item Extend the distinct-seam reflected closure to arbitrary repeated
seams by the completed common-seam transform and finite rational
partial fractions. Keep the generic unreflected cubic reduction open.
\item Add raw harmonic powers with continuous outer shift as a separate
family from elementary-symmetric harmonic numerators. Record the
all-even centered cancellation and the finite-part polynomial theorem.
\item Replace the proposed ``complete Cayley symmetry and retry the
same bounded rows'' step by the proved nonmembership theorem. Keep
$S_6$ and the revised $S_8$ explicitly conjectural.
\item Add a version-specific note for Tin's equation (22), with
\eqref{hp:eq:tin-all-log} as the proof of the corrected sign. Do not
silently change the Laurent convention to conceal the discrepancy.
\end{enumerate}

% END sections/06_audit.tex


% BEGIN sections/07_research.tex
\section{Further research questions and a programme of exact identities}
\label{research:section}

These questions begin at the boundaries of the proved results. They
ask for exact formulas, functional relations, or finite certificates;
none is a conjectural assertion presented as a theorem. They are ordered
roughly by proximity to computations already available in this package.

\subsection{Ordered jets, derivatives, and antiderivatives}

\paragraph{Q1. Differential closure of the depth-four coordinates.}
Equation~\eqref{oq:derivative} already removes $T$ in favor of
$\eta'(a)$, ordinary Hurwitz jets, and the convergent value
$Z_2(2,1;a)$. Derive analogous argument-derivative formulas for
$\rho$ and $\kappa$, then use stuffle and partial fraction identities
to reduce the resulting depth-three and mixed-weight sums. A useful
first result would be a finite differential system with all additive
normalizations fixed by the shift laws. Determine the smallest
\emph{proved spanning} system before asking for arithmetic minimality.

\paragraph{Q2. Two or more blocks of equal slopes at arbitrary depth.}
The all-equal and one-distinguished-slope families have Gamma
generators. At depth four, the fully independent formula identifies
the first coefficients that such generators miss. Seek an exact
two-block generator for slopes
$(p,\ldots,p,q,\ldots,q)$, with the block lengths arbitrary. The target
is a finite collection of lower-depth generating germs independent of
the block lengths, or a proof that a proposed finite collection fails.
The depth-four normal form supplies a mandatory first consistency
specialization.

\paragraph{Q3. Depth five with an explicit coordinate ledger.}
At depth five the finite part needs the cubic jet of $R_2$, the
quadratic jet of $R_3$, and the linear jet of $R_4$. Compute the exact
regularized stuffle constraints on these jets and construct an analogue
of the three-direction reconstruction. The deliverable should include
a convergent representative for every retained coordinate, its shift
recurrence, and a symbolic certificate for each claimed cancellation.
This is a finite linear problem after the analytic regular germs have
been fixed; it is not automatically a period-independence problem.

\paragraph{Q4. Rational-shift character combinations.}
For conductors $2,3,4,6$, apply the exact root-of-unity filter to the
convergent representatives of $\eta,\rho,\kappa,T$. Determine which
character-weighted combinations eliminate the residual ordered
coordinates and reduce to ordinary Hurwitz or Dirichlet $L$ jets.
The shift identities alone give translations, not distribution.
A successful theorem must show the cancellation algebraically and
retain every logarithmic scale term.

\paragraph{Q5. Values on a polar hyperplane.}
The transverse-ray formula assumes all prefix slopes are nonzero.
For a family lying in one prefix hyperplane, compare three explicit
prescriptions: take a finite part in a normal variable first, remove
the ordered poles through $R_d$ first, or approach along a specified
two-parameter surface. Compute their exact differences at depth four
and determine which prescriptions preserve a stated stuffle identity.
No value is defined by substituting a zero denominator into the ray
formula.

\subsection{Coincident products and Gamma primitives}

\paragraph{Q6. The first two-unreflected Stieltjes closure problem.}
Reflection evaluates $\FP\int\psi^2K^{2h+1}$ in this article, while
one arbitrary Stieltjes derivative times a reflected polynomial always
closes. Classify the combinations of
\[
 \FP\int_0^1\gamma_m^{(p)}(x)\gamma_n^{(q)}(x)K(x)^r\,dx
\]
whose completed higher-depth coordinate cancels. Begin with
$m+n\le2$ and total argument derivative order at most two. The output
should be an exact kernel of a finite coordinate map, with a proof of
each cancellation and fixed endpoint coordinates.

\paragraph{Q7. Evaluate the generic cubic Tornheim coordinate.}
The cubic difference \eqref{ct:cubic-difference} reduces a mixed
moment without determining the common unreflected cubic value. Seek a
functional relation for the diagonal Tornheim third derivative in the
incoming cubic formula, or a normalized Barnes-type representation
whose derivative can be proved to equal it. This must be a new analytic
relation, not another rearrangement of an equivalent cubic identity.
The existing Mellin representation is a viable starting point.

\paragraph{Q8. A reflected generating function beyond polynomials.}
Polynomial closure is finite because the cotangent derivative
polynomials form a triangular basis. Determine a nonpolynomial family
$P_\lambda(K)$ for which the completed integral has a holomorphic
parameter generator and controlled coefficientwise endpoint completion.
Trigonometric beta integrals suggest rational or complex-power
families, but their changing endpoint exponents introduce new
resonances. A useful theorem would identify the full local subtraction
before any expansion in $\lambda$.

\paragraph{Q9. Multiple Gamma primitives with fixed constants.}
The log-Gamma moments here are first argument primitives of digamma
moments. Construct analogous Barnes-$G$ or higher-Gamma identities
for repeated primitives, with the integration constant fixed at an
explicit positive point and compatibility checked against the
functional equation. The endpoint constant terms must be included
at every integration by parts. A first target is a closed generator
for the even cotangent moments at the next primitive level.

\subsection{Harmonic powers and spectral regular coefficients}

\paragraph{Q10. Compact all-degree formulas for centered values.}
The finite word recursion proves that every
$F_p(-2r,1/2)$ belongs to the ordinary multiple-zeta algebra.
Seek a bivariate generator in $p$ and $r$ that improves on enumerating
all compositions of $p$. Even a closed formula for the entire sequence
$F_p(-2,1/2)$ would extend the explicit table. State the exact MZV
coordinate space used at each weight rather than assume every value
reduces to products of ordinary zeta values.

\paragraph{Q11. Generalized harmonic products and their regular shifts.}
For a fixed polynomial in $H_n^{(r)}$, $r\ge1$, the same Mellin
calculus produces polynomial finite parts. Which numerator combinations
make every nonpositive even principal part vanish at $a=1/2$?
The raw-power parity proof does not apply automatically when
$H_n^{(r)}$ with $r>1$ is present. Compute the finite Bernoulli
conditions at successive orders and seek a necessary and sufficient
generating criterion. This gives exact cancellation identities rather
than inequalities.

\paragraph{Q12. Cancel the first global Mellin remainder.}
The first regular spectral coefficient at $s=-m$ contains global
information absent from its finite part. Find finite combinations of
harmonic-power series whose remainder integrals cancel by an exact
kernel identity. Proposition~\ref{hp:prop:tin-all-log} is a model:
one differential combination has a normally convergent series and an
explicit relation among higher harmonic Stieltjes constants. Generalize
this cancellation to nonpositive centers and generalized harmonic
numerators, with the normalization checked independently.

\subsection{The Gaussian period candidates}

\paragraph{Q13. Identify additional information detected by the separator.}
The inherited functional survives completion of the entire Cayley
ideal. A useful next computation is to generate selected additional
double-shuffle or distribution relations outside the frozen family,
project them by $\Pi_\uparrow$, and identify a row on which the
functional is nonzero. Such a row is a concrete source of information
missing from the present search; it is not by itself a proof of $S_6$.
Continue until either an exact rational proof certificate for the
target is obtained or a new fully specified separating certificate is
proved. Higher-depth intermediate rows are permitted, but need not be
assumed necessary a priori.

\paragraph{Q14. A functional bridge to $S_6$ and the revised $S_8$.}
Seek a parameter-dependent polylogarithm identity, differential system
with enough exact boundary data, or regulator relation whose
specialization is the short candidate. Independently, repeat the
precisely defined formal quotient analysis for the revised weight-nine
target. The weight-seven certificate cannot be extrapolated to it.
Any new numerical relation remains a candidate until such a functional
or exact symbolic proof is supplied.

\paragraph{Q15. Formal verification of the finite-part and ideal mechanisms.}
The local completion lemma and the depth-preserving ideal-intersection
lemma are reusable proof units. Formalize those units in a proof
assistant, then import the finite rational certificate with an exact
integer parser. For the analytic identities, formalize normal
convergence after a finite local subtraction before importing
coefficient formulas. This would turn the current reproducible
computations into a smaller formally checked core without confusing
floating-point diagnostics with theorem proofs.

% END sections/07_research.tex


% BEGIN sections/08_verification.tex
\section{Proof dependencies and reproducible verification}
\label{verification:section}

The all-index analytic statements are proved by finite local
subtraction, normal convergence, and exact coefficient extraction.
The formal ideal statement combines an all-weight algebraic argument
with a finite exact rational certificate. The additional computations
have two purposes: to detect normalization or transcription mistakes,
and to make the finite algebraic part independently replayable.
They do not replace the analytic arguments.

\subsection{Exact calculations}

The ordered verifier solves the regularized stuffle equations for the
required jets, extracts all Laurent coefficients through the finite
part, checks the diagonal Gamma specialization and shift identities,
and reconstructs the three-direction matrix. It passes 23 exact
symbolic assertions, including a control that rejects an intentionally
changed coefficient of $\kappa$.

The harmonic symbolic verifier independently constructs the low-weight
multiple-polylogarithm expansions from their differential equations
and fixed endpoint constants. It checks the Bernoulli principal
coefficients, the displayed finite parts, parameter transport, and
selected higher centered cancellations. It passes 171 exact symbolic
assertions. The all-index parity proof remains the reason that every
centered even pole vanishes.

The reflected verifier checks 15 exact derivative-polynomial identities
in addition to its numerical comparisons. The complete Cayley verifier
regenerates all 5,131 inherited rows, verifies all 2,546 basis
differences in the ideal intersection, and obtains the exact nonzero
target pairing in Section~\ref{gf:sec}. It also checks 780 independent
Eulerian/convention cases and 2,150 shuffle products. Both retained
upstream code modules are byte-for-byte copies of their pinned Git
blobs. The verifier checks their SHA-256 hashes and the exact
functional data before carrying out the calculation.

\subsection{Independent numerical diagnostics}

\begin{center}
\begin{tabular}{@{}p{0.20\textwidth}p{0.45\textwidth}p{0.24\textwidth}@{}}
\toprule
Family & Independent evaluation & Largest observed discrepancy\\
\midrule
Ordered depth four & Original ordered sums with Euler--Maclaurin tails;
four directional finite parts and four shift checks &
$6.54\cdot10^{-34}$\\[5pt]
Reflected moments & Central ordinary quadrature plus independent
endpoint Laurent-log integration; 30 checks &
$4.38\cdot10^{-47}$\\[5pt]
Harmonic powers & Original Dirichlet head and digamma asymptotic tail;
72 Laurent coefficients including larger-cutoff replays &
$3.41\cdot10^{-41}$\\
\bottomrule
\end{tabular}
\end{center}

The ordered run uses 55 working decimal digits with 35 guard digits
for its local coefficient calculation, and tests shifts $a=1,2$.
The reflected run uses 50 decimal digits and includes higher
Stieltjes/argument derivative pairs that test the nonconstant resonant
amplitude. The harmonic run uses 58 decimal digits, shifts $1/2$ and
$3/4$, and spectral centers $0,-1,-2$; its independent replay increases
the finite head, tail order, and Cauchy sample count. The attached JSON
files record all numerical settings and individual residuals.

These are floating-point diagnostics, not interval enclosures. The
numbers in the last column are observed discrepancies, not rigorous
error bounds. The finite symbolic checks are also not proof-assistant
formalizations of the analytic theorems. Separate derivations and
cross-reviews checked the ordered jet equations, the harmonic local
table, their factorial conventions, and their shift identities.

\subsection{Files and replay commands}

The package includes the modular article sources, a self-contained
\src{Exact_Resonant_Identities.tex}, the compiled PDF, exact certificate
data, verification scripts and receipts, a source manifest, a claims
ledger, and integration notes. The self-contained source inlines the
section and bibliography inputs, so its compilation does not require
network retrieval or access to the inspected repository.

The numerical and symbolic scripts use Python with \src{mpmath} and
\src{sympy}; the complete Cayley verifier requires only the Python
standard library. From the package root:
\begin{verbatim}
python3 build.py
python3 verification/verify_all.py --exact-only
python3 verification/verify_all.py
\end{verbatim}
The last command includes the higher-index reflected checks and the
full independent numerical runs. The detailed README lists separate
commands so a reader can replay only the result under review. No search
matrix, modular rank cache, numerical polylogarithm conjecture check,
or external service is required for the exact Cayley proof.

The top-level checksums describe the delivered files. Verification
receipts contain elapsed times and may change when replayed; source
integrity checks for the actual certificate are independent of those
timings. No numerical threshold is used to decide the nonzero exact
integer pairing.

% END sections/08_verification.tex

\clearpage

% BEGIN references.tex
\begin{thebibliography}{99}
\addcontentsline{toc}{section}{References}

\bibitem{RepoMain}
V. Reshetnikov and the ProveIt research programme,
\emph{Polylogarithms and their Arithmetic Bridges}, canonical manuscript,
status records, and integrated research chapters. Repository snapshot
\texttt{5a790187c8e186e41e2b990b4941cb7a1a3c7b6b}, 11 October 2026.
\href{https://github.com/VladimirReshetnikov/ProveIt/tree/5a790187c8e186e41e2b990b4941cb7a1a3c7b6b/Analysis/Polylogarithms/docs/manuscript}{Pinned manuscript directory}.

\bibitem{RepoOrdered}
ProveIt research continuation,
\emph{Ordered Hurwitz Resonances and Cubic Stieltjes Products:
Exact Laurent coefficients, Gamma generating functions, sharp curve
dependence, and a polylogarithmic Mellin bridge}, 10 October 2026.
Incoming archive \src{ProveIt_Ordered_Hurwitz_Resonances_2026-10-10.zip};
particularly \src{sections/02_ordered.tex} and
\src{sections/07_questions.tex}, Questions 1--2.
\href{https://github.com/VladimirReshetnikov/ProveIt/blob/5a790187c8e186e41e2b990b4941cb7a1a3c7b6b/docs/incoming/ProveIt_Ordered_Hurwitz_Resonances_2026-10-10.zip}{Pinned archive}.

\bibitem{RepoTriple}
ProveIt research continuation,
\emph{Triple Stieltjes Correlations: Colored Tornheim Jets, Harmonic
Contact Terms, and Exact Gamma Reductions}, 10 October 2026.
Incoming archive \src{ProveIt_Triple_Stieltjes_Tornheim_2026-10-10.zip};
particularly \src{sections/08-reflection.tex} and the cubic and
normalization sections.
\href{https://github.com/VladimirReshetnikov/ProveIt/blob/5a790187c8e186e41e2b990b4941cb7a1a3c7b6b/docs/incoming/ProveIt_Triple_Stieltjes_Tornheim_2026-10-10.zip}{Pinned archive}.

\bibitem{Saha}
B. Saha, Multiple Stieltjes constants and Laurent type expansion of the
multiple zeta functions at integer points,
\emph{Selecta Math. (N.S.)} \textbf{28} (2022), article 6.
\url{https://doi.org/10.1007/s00029-021-00719-1}.
Author preprint: \url{https://arxiv.org/abs/1902.04389}.

\bibitem{MOW}
K. Matsumoto, T. Onozuka, and I. Wakabayashi,
Laurent series expansions of multiple zeta-functions of Euler--Zagier
type at integer points, \emph{Math. Z.} \textbf{295} (2020), 623--642.
\url{https://doi.org/10.1007/s00209-019-02337-2}.
Author preprint: \url{https://arxiv.org/abs/1601.05918}.

\bibitem{EspinosaMoll1}
O. Espinosa and V. H. Moll,
On some integrals involving the Hurwitz zeta function: Part 1,
\emph{Ramanujan J.} \textbf{6} (2002), 159--188.
\url{https://www.math.tulane.edu/~vhm/papers_html/hurwitz1.pdf}.

\bibitem{EspinosaMoll2}
O. R. Espinosa and V. H. Moll,
On some integrals involving the Hurwitz zeta function: Part 2,
\emph{Ramanujan J.} \textbf{6} (2002), 449--468.
Author preprint: \url{https://arxiv.org/abs/math/0107082}.

\bibitem{Kargin}
L. Karg\i n, A. Dil, M. Cenkci, and M. Can,
On the Stieltjes constants with respect to harmonic zeta functions,
\emph{J. Math. Anal. Appl.} \textbf{525} (2023), article 127302.
\url{https://doi.org/10.1016/j.jmaa.2023.127302}.
Author preprint: \url{https://arxiv.org/abs/2304.03517}.


% BEGIN references_harmonic.tex
% Bibitems to merge into the parent article's bibliography.
\bibitem{hp:Tin2025}
Lo Ho Tin,
\emph{On sums involving powers of harmonic numbers},
arXiv:2507.03058v1, 2025.
\url{https://arxiv.org/abs/2507.03058v1}.
The sign comparison in Section~\ref{hp:sec:tin-audit} concerns
version 1, equation (22), printed page 5.

\bibitem{hp:IKZ2006}
Kentaro Ihara, Masanobu Kaneko, and Don Zagier,
\emph{Derivation and double shuffle relations for multiple zeta values},
Compositio Mathematica \textbf{142} (2006), no.~2, 307--338.
\url{https://doi.org/10.1112/S0010437X0500182X}.

\bibitem{hp:DLMF}
NIST Digital Library of Mathematical Functions,
\emph{Chapter 5: Gamma Function}, \S5.11, Asymptotic Expansions.
\url{https://dlmf.nist.gov/5.11}.

% END references_harmonic.tex


% BEGIN references_cayley.tex
\bibitem{gf:Bounded}
ProveIt research contribution,
\emph{A certified obstruction in the weight-seven reduction search},
in the extremal cyclic-descent continuation,
\texttt{article/s6\_obstruction.tex}.
Pinned repository commit
\texttt{5a790187c8e186e41e2b990b4941cb7a1a3c7b6b}.
\href{https://github.com/VladimirReshetnikov/ProveIt/blob/5a790187c8e186e41e2b990b4941cb7a1a3c7b6b/Analysis/Polylogarithms/docs/reports/formal-reductions-continuation/continuations/extremal-cyclic-descent/article/s6_obstruction.tex}{Source article}.

\bibitem{gf:Depth}
ProveIt research contribution,
\emph{A Cayley normal form that preserves series depth},
in the polylogarithm--Stieltjes critical-transition continuation,
\texttt{article/05\_cayley\_depth.tex}.
Pinned repository commit
\texttt{5a790187c8e186e41e2b990b4941cb7a1a3c7b6b}.
\href{https://github.com/VladimirReshetnikov/ProveIt/blob/5a790187c8e186e41e2b990b4941cb7a1a3c7b6b/Analysis/Polylogarithms/docs/reports/uniform-transition-continuation/continuations/polylog-stieltjes-critical-transition/article/05_cayley_depth.tex}{Source article}.

\bibitem{gf:Radford}
D.~E.~Radford,
A natural ring basis for the shuffle algebra and an application to group
schemes, \emph{Journal of Algebra} \textbf{58} (1979), no.~2, 432--454.
\href{https://doi.org/10.1016/0021-8693(79)90171-6}{doi:10.1016/0021-8693(79)90171-6}.

\bibitem{gf:Patras}
F.~Patras,
La d\'ecomposition en poids des alg\`ebres de Hopf,
\emph{Annales de l'Institut Fourier} \textbf{43} (1993), no.~4, 1067--1087.
\href{https://doi.org/10.5802/aif.1365}{doi:10.5802/aif.1365}.


% END references_cayley.tex


\bibitem{DLMFHurwitz}
NIST Digital Library of Mathematical Functions,
\S25.11, Hurwitz zeta function.
\url{https://dlmf.nist.gov/25.11}. Accessed 11 October 2026.

\bibitem{DLMFPolyGamma}
NIST Digital Library of Mathematical Functions,
\S5.15, Polygamma functions.
\url{https://dlmf.nist.gov/5.15}. Accessed 11 October 2026.

\bibitem{DLMFGamma}
NIST Digital Library of Mathematical Functions,
\S5.5, Functional relations for the Gamma function.
\url{https://dlmf.nist.gov/5.5}. Accessed 11 October 2026.

\bibitem{DLMFBeta}
NIST Digital Library of Mathematical Functions,
\S5.12, Beta function.
\url{https://dlmf.nist.gov/5.12}. Accessed 11 October 2026.

\bibitem{DLMFZetaReflection}
NIST Digital Library of Mathematical Functions,
\S25.4, Riemann zeta reflection formulas.
\url{https://dlmf.nist.gov/25.4}. Accessed 11 October 2026.
\end{thebibliography}

% END references.tex

\end{document}
